%%%%%%%% ICML 2025 EXAMPLE LATEX SUBMISSION FILE %%%%%%%%%%%%%%%%%

\documentclass{article}

% Recommended, but optional, packages for figures and better typesetting:
\usepackage{microtype}
\usepackage{graphicx}
% \usepackage{subfigure}
\usepackage{booktabs} % for professional tables

% hyperref makes hyperlinks in the resulting PDF.
% If your build breaks (sometimes temporarily if a hyperlink spans a page)
% please comment out the following usepackage line and replace
% \usepackage{icml2025} with \usepackage[nohyperref]{icml2025} above.
\usepackage{hyperref}


% Attempt to make hyperref and algorithmic work together better:
\newcommand{\theHalgorithm}{\arabic{algorithm}}

% Use the following line for the initial blind version submitted for review:
\usepackage{icml2025}

% If accepted, instead use the following line for the camera-ready submission:
% \usepackage[accepted]{icml2025}

% For theorems and such
\usepackage{amsmath}
\usepackage{amssymb}
\usepackage{mathtools}
\usepackage{amsthm}

% if you use cleveref..
\usepackage[capitalize,noabbrev]{cleveref}

%%%%%%%%%%%%%%%%%%%%%%%%%%%%%%%%
% THEOREMS
%%%%%%%%%%%%%%%%%%%%%%%%%%%%%%%%
\theoremstyle{plain}
\newtheorem{theorem}{Theorem}[section]
\newtheorem{proposition}[theorem]{Proposition}
\newtheorem{lemma}[theorem]{Lemma}
\newtheorem{corollary}[theorem]{Corollary}
\theoremstyle{definition}
\newtheorem{definition}[theorem]{Definition}
\newtheorem{assumption}[theorem]{Assumption}
\theoremstyle{remark}
\newtheorem{remark}[theorem]{Remark}

% Todonotes is useful during development; simply uncomment the next line
%    and comment out the line below the next line to turn off comments
%\usepackage[disable,textsize=tiny]{todonotes}
\usepackage[textsize=tiny]{todonotes}

%new
\usepackage{xspace}
% \usepackage{graphicx}
\usepackage{subfig}
% 删除了本来的% \usepackage{subfigure}
\newcommand{\consensustester}{Agora\xspace} 
\newcommand{\agenta}{Orchestrator\xspace} 
\newcommand{\agentb}{Strategy\xspace} 
\newcommand{\agentc}{TestGen\xspace} 
\newcommand{\bfagenta}{\textbf{Orchestrator}\xspace} 
\newcommand{\bfagentb}{\textbf{Strategy}\xspace} 
\newcommand{\bfagentc}{\textbf{TestGen}\xspace} 
\usepackage{tikz}
\newcommand{\circled}[1]{\tikz[baseline=(char.base)]{
  \node[shape=circle,draw,inner sep=0.2pt,minimum size=0.8em, font=\scriptsize] (char) {#1};
}}
% \newcommand{\circledsmall}[1]{%
% \tikz[baseline=(char.base)]{
%   \node[circle,draw,inner sep=0pt,font=\scriptsize] (char) {#1};
% }}
\newcommand{\BugAnalysisBox}[3]{%
  \begin{tcolorbox}[
    colframe=orange!60!black,    % 边框：深橙
    coltitle=orange!80!black,     % 标题文字：橙黑混合
    colbacktitle=orange!20!white, % 标题栏背景：更淡的橙
    colback=white, % 设置主体背景为纯白
    title={#1}, 
    fonttitle=\bfseries,
    nobeforeafter, center title,
    enhanced,breakable
  ]
    \textbf{Analysis}: #2 
  \end{tcolorbox}
}
\usepackage[most]{tcolorbox}

\newcommand{\PatternBox}[3]{%
  \begin{tcolorbox}[
    enhanced,
    breakable,
    colframe=blue!50!black,    % 边框颜色：深蓝
    colback=white,             % 主体背景：纯白
    coltitle=white,            % 标题文字：白色
    colbacktitle=blue!50!black, % 标题背景：深蓝（与边框一致）
    title={#1}, 
    fonttitle=\bfseries\sffamily, % 标题字体：粗体无衬线
    sharp corners,             % 方角（图片中看起来是直角）
    boxrule=0.8pt,             % 边框线条粗细
    titlerule=0pt,             % 去掉标题与主体之间的分割线，使其浑然一体
    center title,              % 标题居中
    left=10pt, right=10pt, top=10pt, bottom=10pt % 内部留白
  ]
    #2 % 这里填入内容
  \end{tcolorbox}
}

%Orchestrator Agent bug-exploitation state-analyzer
%Strategy Agent Constraint Scenarios
%TestGen Agent Loop
%实验需要braft zookeeper
%三个级别的bug
% \usepackage{etoc}
\usepackage[table]{xcolor}

\usepackage[most]{tcolorbox} % 核心
\usepackage{xcolor}          % 颜色基础
\usepackage{fancyvrb} % 提供 Verbatim 环境
% 补充缺失的颜色定义
\definecolor{SeaCornSoft}{HTML}{FFF8DC} 
\definecolor{RoyalAquaDeep}{HTML}{006064}
\newtcolorbox{agentbox}[1][]{
  enhanced, breakable, drop shadow,
  width=\linewidth, boxrule=0.8pt,
  colback=SeaCornSoft,       % 背景
  colframe=RoyalAquaDeep,    % 边框
  left=4pt,right=4pt,top=4pt,bottom=4pt,
  fonttitle=\small,
  title={\bfseries #1}
}
%new


% The \icmltitle you define below is probably too long as a header.
% Therefore, a short form for the running title is supplied here:
\icmltitlerunning{\consensustester: Toward Autonomous Bug Detection in Production-Grade Consensus Protocols with LLM Agents}

\begin{document}

\twocolumn[
% \icmltitle{\consensustester: Toward Autonomous Bug Detection in Production-Grade Consensus Protocols with LLM Agents}
\icmltitle{
    \begin{tabular}{cc}
        \raisebox{-15pt}{\includegraphics[height=35pt]{pic/agora.png}} & 
        \begin{tabular}{c}
             \consensustester: Toward Autonomous Bug Detection in 
            Production-Grade\\ Consensus Protocols with LLM Agents
        \end{tabular}
    \end{tabular}
}
% \icmltitle{
%     \begin{center}
%         \begin{minipage}{0.1\textwidth}
%             \centering
%             \includegraphics[height=40pt]{pic/agora.png}
%         \end{minipage}%
%         \begin{minipage}{0.9\textwidth}
%             \centering
%             \consensustester: Toward Autonomous Bug Detection in Production-Grade Consensus Protocols with LLM Agents
%         \end{minipage}
%     \end{center}

% }

% It is OKAY to include author information, even for blind
% submissions: the style file will automatically remove it for you
% unless you've provided the [accepted] option to the icml2025
% package.

% List of affiliations: The first argument should be a (short)
% identifier you will use later to specify author affiliations
% Academic affiliations should list Department, University, City, Region, Country
% Industry affiliations should list Company, City, Region, Country

% You can specify symbols, otherwise they are numbered in order.
% Ideally, you should not use this facility. Affiliations will be numbered
% in order of appearance and this is the preferred way.
\icmlsetsymbol{equal}{*}

\begin{icmlauthorlist}
\icmlauthor{Firstname1 Lastname1}{equal,yyy}
\icmlauthor{Firstname2 Lastname2}{equal,yyy,comp}
\icmlauthor{Firstname3 Lastname3}{comp}
\icmlauthor{Firstname4 Lastname4}{sch}
\icmlauthor{Firstname5 Lastname5}{yyy}
\icmlauthor{Firstname6 Lastname6}{sch,yyy,comp}
\icmlauthor{Firstname7 Lastname7}{comp}
%\icmlauthor{}{sch}
\icmlauthor{Firstname8 Lastname8}{sch}
\icmlauthor{Firstname8 Lastname8}{yyy,comp}
%\icmlauthor{}{sch}
%\icmlauthor{}{sch}
\end{icmlauthorlist}

\icmlaffiliation{yyy}{Department of XXX, University of YYY, Location, Country}
\icmlaffiliation{comp}{Company Name, Location, Country}
\icmlaffiliation{sch}{School of ZZZ, Institute of WWW, Location, Country}

\icmlcorrespondingauthor{Firstname1 Lastname1}{first1.last1@xxx.edu}
\icmlcorrespondingauthor{Firstname2 Lastname2}{first2.last2@www.uk}

% You may provide any keywords that you
% find helpful for describing your paper; these are used to populate
% the "keywords" metadata in the PDF but will not be shown in the document
\icmlkeywords{Machine Learning, ICML}

\vskip 0.3in
]

% this must go after the closing bracket ] following \twocolumn[ ...

% This command actually creates the footnote in the first column
% listing the affiliations and the copyright notice.
% The command takes one argument, which is text to display at the start of the footnote.
% The \icmlEqualContribution command is standard text for equal contribution.
% Remove it (just {}) if you do not need this facility.

\printAffiliationsAndNotice{}  % leave blank if no need to mention equal contribution
% \printAffiliationsAndNotice{\icmlEqualContribution} % otherwise use the standard text.




% 模仿 AlphaFold + KNighter 的写法
%   - AlphaFold: 开篇格局大，强调"X年未解难题"，结尾升华
%   - KNighter: 指出LLM直接应用的问题，提出"合成"范式，强调闭环精炼
%
% Abstract 五段式结构：
%   1. The Grand Challenge - 背景与痛点（格局要大）
%   2. The Gap - 现有方法的局限（形式化太费力 与Fuzzing缺语义）
%   3. The Insight - 核心洞察（LLM直接用有问题 → 测试合成范式）
%   4. The Methodology - 方法论（闭环精炼流水线）
%   5. The Results - 实验结果（X个bug）
%
% 关键卖点：
%   - "Test Synthesis" 范式（不是直接用LLM分析，而是合成测试Agent）
%   - "Closed-loop Refinement" 闭环精炼（Strategy → TestGen → Feedback → 迭代）
%   - "Democratize" 民主化（降低分布式系统验证的门槛）


% % 首先是背景与痛点
% % 模仿AlphaFold开篇，把问题上升到"现代基础设施基石"的高度
% Distributed consensus protocols (e.g., Raft, PBFT) serve as the backbone of modern cloud infrastructure, enabling reliable coordination across databases, blockchain networks, and mission-critical services that billions of users depend on daily. 
% However, verifying their correctness remains a grand challenge, bottlenecked by the state space explosion that makes exhaustive testing computationally intractable.

%多少天

% % 现有方法的局限
% % 形成对比——形式化验证太费力 与 Fuzzing缺语义理解
% % 关键词: fundamental dilemma, prohibitive manual effort, semantic understanding
% Existing techniques face a fundamental dilemma: formal verification (e.g., TLA+, model checking) requires prohibitive manual effort---often months per system---while coverage-based fuzzing lacks the semantic understanding to navigate complex protocol states, missing deep logic bugs that only manifest under specific interleavings of network partitions, message reorderings, and crash-recovery sequences.

% % 核心问题
% % 指出LLM直接应用的问题（上下文限制、幻觉），提出"测试合成"范式
% % 这是区别于直接用LLM的关键insight，模仿KNighter的"合成checker"思路
% While Large Language Models (LLMs) offer powerful semantic reasoning capabilities, directly applying them to protocol testing is hindered by context limitations and hallucination problems. 
% We present Consensus-Tester, which shifts the paradigm from direct LLM-based analysis to \textbf{automated test synthesis}---leveraging LLMs to generate protocol-aware adversarial test agents rather than analyzing code directly.

% % 方法论
% % 写作思路: 强调"闭环精炼"是核心机制，描述pipeline
% Consensus-Tester leverages the latent protocol priors within LLMs to synthesize executable, adversarial test cases. Central to our approach is a \textbf{closed-loop refinement pipeline}: a Strategy Agent generates semantic-aware attack scenarios by reasoning about protocol invariants, a TestGen Agent synthesizes executable tests with self-healing capabilities, and execution feedback iteratively refines the testing strategy. This closed loop guides the system toward rare corner cases at the ``crash boundary'' of protocol implementations.

% % 实验结果
% We evaluate Consensus-Tester across 10+ widely-deployed consensus implementations spanning both crash fault-tolerant (CFT) protocols (hashicorp-raft, etcd/raft, raft-rs, ZooKeeper) and Byzantine fault-tolerant (BFT) protocols (CometBFT, HotStuff, BFT-SMaRt). Consensus-Tester has discovered \textbf{9} bugs, with \textbf{2} confirmed. These bugs span diverse failure modes: safety violations including split-brain after crash recovery, state machine divergence, and quorum certificate replay attacks; liveness bugs such as deadlocks; and silent failures where errors are not propagated to callers. 
% Our results establish a new paradigm: AI-driven protocol testing that democratizes access to rigorous verification, transforming a task that once required deep distributed systems expertise into an automated, scalable process.
\begin{abstract}
% \bfagenta
% \bfagentb
% \bfagentc
% \agenta
% \agentb
% \agentc

%如果没找到bug，那么如果调整

% While LLM-based agents have shown promise in code analysis and bug detection, existing approaches often focus on superficial issues and lack structured reasoning over long-horizon, protocol-level behaviors. Detecting deep, state dependent bugs in complex software systems remains a challenging reasoning problem for current Large Language Models (LLMs), especially in domains governed by consensus protocol semantics. 
% To mitigate this gap, we present \consensustester, a domain-aware multi-agent testing framework that integrates hypothesis-driven testing with LLM-based agents to enable systematic reasoning over consensus protocols. 
% \consensustester decomposes the bug detection process into three specialized agents responsible for workflow control (\bfagenta), consensus scenarios properties (\bfagentb), and test generation (\bfagentc). Through explicit role separation and agent coordination, the framework enables structured exploration of complex protocol states and behaviors beyond local code patterns.
% We evaluate \consensustester using multiple LLMs on four widely used consensus protocol implementations written in different programming languages. The framework uncovers 15 previously unknown zero-day bugs and outperforms baselines. Ablation studies further demonstrate the importance of agent decomposition and domain knowledge integration for discovering deep protocol-level bugs.

Consensus protocols form the backbone of distributed systems and blockchains, where implementation bugs can cause data corruption and financial losses. 
While LLM-based approaches show promise in code analysis, they struggle 
with deep protocol-level logic bugs involving complex state-dependent behaviors across multiple execution stages.
We present \consensustester, a domain-aware multi-agent framework that integrates hypothesis-driven testing with LLM capabilities for systematic protocol verification. \consensustester employs specialized agents that collaboratively explore protocol state spaces, synthesize attack scenarios using domain-specific constraints, and validate findings through iterative refinement. This explicit role separation enables reasoning about global protocol invariants beyond single-function code analysis.
We evaluate \consensustester on four consensus implementations (Raft, EPaxos, HotStuff, BullShark) using four state-of-the-art LLMs. \consensustester discovers 15 previously unknown protocol-level logic bugs that violate safety properties, while existing LLM-based agents fail to detect any such protocol-level logic bugs. Our results demonstrate that domain-aware multi-agent collaboration is essential for detecting deep logic bugs in complex protocols.
\end{abstract}



\section{Introduction}
\label{sec:intro}
%sound complete 不写！
%fuzzing写不写
%static? nonono
%safety liveness?
%bug分成三个level
%etcd的livenss要加上我的场景！
%hotstuff 2 3不是bug
%epaxos 7不是bug
% 强调去掉了一些components，我们的也能发现low-level bugs！！！！对比baselines
%bug不准 还得重新迭代确认
% 图：https://docs.google.com/presentation/d/1nAUFfJY8hSDMMQwjzQPGgWLzvyxqDDDByO5SGkVcEO4/edit?usp=sharing
% https://docs.google.com/presentation/d/1uWBcTgTzSNc9pyRHX0qnJanlKL-MzyXT8KjIY4LDkjk/edit?usp=sharing

%不是MCP 是function call
%不足的地方：bug分析，token分析，bug有些是错的，digest,ablation不是特别有效
%===================

%第一段 应该是共识多么重要，BFT CFT

%讲一下共识protocols是整个分布式系统的基础，其中 crash fault-tolerant (CFT) protocols，可以支持分布式数据库，存储系统和集群协调与配置管理，保证了线上交易 数据库存储的正确性,而Byzantine fault-tolerant (BFT) 则是blockchain system，设计区块链交易的安全。
%因此，当共识协议和其实现出现了问题，可能会导致严重的问题

%目前的方法为何不行
%虽然进行分布式协议代码的找bug一直都是研究的重点。但是，然而由于分布式系统实现的复杂情况以及共识协议本身的复杂，导致bug很难被发现。
%用了很久的库很多年的库etcd 即便很多需要很多年，

%LLM的方法和agent 
%balance


Consensus protocols, including Crash Fault-Tolerant (CFT)~\cite{Lamport1998ThePP,Lamport2001PaxosMS,Moraru2013ThereIM,Ongaro2014InSO} and Byzantine Fault-Tolerant (BFT)~\cite{Castro1999PracticalBF, Lamport2011ByzantizingPB, Yin2018HotStuffBC,Danezis2021NarwhalAT} protocols, serve as the cornerstone of modern distributed systems, ensuring consistency and availability in the presence of faults. Among these, the CFT protocols are widely adopted in distributed databases~\cite{taft2020cockroachdb}, storage systems~\cite{corbett2013spanner}, and cluster coordination services~\cite{hunt2010zookeeper} to guarantee the correctness of online transactions and data storage. Conversely, the BFT protocols form the backbone of blockchain systems~\cite{nakamoto2008bitcoin}, designed to secure transactions against malicious actors. 

%TLA+ 不写了吧 model checking
Consequently, any defects in the implementation of these consensus protocols can lead to severe consequences, ranging from data corruption to significant financial losses~\cite{zheng2017paxosstore}.
Despite the critical importance of these systems, bug detection for distributed protocol implementations remains a formidable challenge and has long been a focal point of research. The inherent complexity of distributed systems, coupled with the intricate logic of consensus algorithms, makes subtle bugs notoriously difficult to detect~\cite{Ongaro2014InSO}. Even widely used and mature libraries (e.g., etcd~\cite{etcdioraft}), which have been under development for many years, are still found to contain critical bugs that require experienced engineers extensive periods to uncover and resolve~\cite{Zhou2021FoundationDBAD}.



\begin{figure*}[!th]
  \centering
  \includegraphics[width=\textwidth]{pic/agora_overview1.pdf}

  
  \caption{\textbf{Left} is the baseline overview. \textbf{Middle} shows the input and result comparison where baselines could only detect \textbf{low-level implementation bugs}, but \consensustester could detect critical \textbf{protocol-level logic bugs}. \textbf{Right} is the overview of \consensustester, which consists of three agents: \agenta agent, \agentb agent and \agentc agent.}
  \label{fig:overview}
  %为啥有<>
\end{figure*}

%balance human knowedge和llm?
%我们得想想问题，agent结合，之前的没有
In recent years, the rapid advancement of Large Language Models (LLMs)~\cite{Achiam2023GPT4TR,DeepSeekAI2025DeepSeekR1IR} and autonomous agent systems~\cite{yao2022react,cheng2024exploring} has revolutionized software engineering tasks, helping engineers save significant time and rapidly perform tasks such as code analysis~\cite{feng2020codebert} and bug detection~\cite{chen2021evaluating}. While LLMs and agents have demonstrated strong capabilities in identifying common coding errors and vulnerabilities~\cite{Kaniewski2025ASL}, their effectiveness remains limited: they primarily excel at detecting superficial or localized bugs, but struggle to reason about the complex, state dependent logic inherent in distributed consensus protocols~\cite{Wang2024SanitizingLL,Wang2024LLMDFAAD}.
Therefore, to fully leverage the general code understanding capabilities of LLMs and efficiently detect bugs in consensus protocols and their implementations, it is necessary to integrate domain knowledge of distributed consensus with multi-agent systems~\cite{Yang2025KernelGPTEK,Yang2025KNighterTS}. Such a fusion enables the identification of subtle errors that may otherwise lead to severe data inconsistencies or substantial economic losses, thereby ensuring system correctness and security for researchers and engineers.




%我们如何做的，分析bug pattern
% Hypothesis-Driven Testing (HDT)~\cite{Kell2004HereIT}
% (1)(2)(3)
% 因此在本文中，我们首先对共识协议的bug进行了分析出了相应的pattern，对共识的bug出现的constrains进行了总结，并且基于testing principle of hypothesis-driven testing (HDT)~\cite{Kell2004HereIT}构建了一个mutil-agent 系统\consensustester 来更好地 融合domian knowledge和LLM强大的code理解能力和泛化能力对共识协议的代码仓库进行分析，来高效的寻找共识协议的bug。具体而言，如图1所示，，我们的\consensustester,由三个agent来实现流程控制，共识特性和测试生成的解耦，进行高效地寻找共识协议深层次的Deep Bug Hunting：
% \textit{(1)}\bfagenta agent负责agents交互，对已有bug总结
% \textit{(2)}\bfagentb agent负责生成特殊的的共识协议的bug场景
% \textit{(3)}\bfagentc agent基于bug场景generate 相应的unitest，测试出bug。
Therefore, in this paper, we first analyze bugs in consensus protocols and identify corresponding patterns, and summarize the constraints under which bugs occur in consensus protocols. Then, based on the testing principle of hypothesis-driven testing (HDT)~\cite{Kell2004HereIT}, we are the first to construct a multi-agent system \consensustester to better integrate domain knowledge with the strong code understanding and generalization capabilities of LLMs to analyze consensus protocol code repositories and efficiently discover bugs in consensus protocols. Specifically, as shown in Figure~\ref{fig:overview}, our \consensustester is composed of three agents to realize the decoupling of workflow control, consensus scenarios properties, and test generation, enabling efficient deep bug hunting in consensus protocols:
\textit{(1)} \bfagenta agent is responsible for agent interactions and summarizing existing bugs, \textit{(2)} \bfagentb agent is responsible for generating special consensus protocol bug scenarios, \textit{(3)} \bfagentc agent generates corresponding unit tests based on the bug scenarios to expose bugs.
%production level
We also used our \consensustester system to conduct extensive experiments on the repositories of four popular and widely used research- and production-level consensus protocols to validate the effectiveness of our approach. Our main contributions are summarized as follows:


\begin{figure*}[!th]
  \centering
  \includegraphics[width=\textwidth]{pic/agora_overview2.pdf}
  \caption{The overview of \consensustester workflow, which consists of three agents and twelve steps/subtasks (following the time order, from Step \circled{1} to Step \circled{12}). \bfagenta agent is responsible for bug exploitation, coordinating the agents, maintaining global state, and sending and receiving messages among agents; \bfagentb agent is responsible for generating attack scenarios based on the analysis results and the information it receives; \bfagentc agent is responsible for generating valid unit tests through a reflection loop to detect bugs in consensus systems. Each agent consists of four subtasks.}%
  %结果是尖括号
  \label{fig:overview2}
\end{figure*}

%contribution
%可解释性
% 多种不同的变成语言  go, Rust

\begin{itemize}
\item We are the first to combine domain knowledge of distributed consensus with LLM-based agents to detect the protocol-level logic bugs.

\item We implement a flexible agent framework, \consensustester, which effectively follows the popular hypothesis-driven testing (HDT) paradigm and employs three agents (\agenta, \agentb, and \agentc) to automatically perform code bug detections for multiple consensus protocol code repositories.

\item We evaluate \consensustester using multiple LLM models and conduct extensive experiments on four popular consensus protocol repositories implemented in different programming languages. Our approach discovers 15 zero-day bugs, demonstrating that \consensustester can significantly improve the safety of distributed systems. 

\item We conduct comparisons with multiple baselines and ablation studies, further validating the effectiveness of our approach and the necessity of each component in integrating LLMs with domain knowledge to uncover deep bugs for consensus protocols. %第四个是不是有点多
\end{itemize}



\section{Preliminary and Related Works}
\label{sec:pre}

\subsection{Hypothesis-Driven Testing}

% 对比传统测试：验证软件"做了它应该做的事" ,Hypothesis-Driven Testing (HDT)可以探索软件在什么条件下会失败,帮助实现更高的code coverage，更加完备的发现bug，去符合  (Vulnerability Hypothesis)：
% % ● 一个漏洞假设 $H = (C, A, E, O)$ 是一个四元组：$C$: Preconditions（前置条件）—— 触发漏洞所需的初始状态
% % ●$A$: Attack Actions（攻击动作）—— 触发漏洞的操作序列
% % ●$E$: Expected Buggy Behavior（预期错误行为）—— 如果漏洞存在的可观察症状
% % ●$O$: Oracle Assertions（断言）—— 验证漏洞的判定条件

% 因此构建一个loop符合相应的HDT的范式不断地自动发现bug，是我们\consensustester的目标

Compared with traditional testing, which only focuses on verifying software does what it is supposed to do, hypothesis-driven testing (HDT)~\cite{Kell2004HereIT} aims to explore under what conditions the software may fail. By explicitly formulating vulnerability hypotheses, HDT enables higher code coverage and more systematic bug discovery. A vulnerability hypothesis is defined as a tuple $H = (C, A, E, O)$, where $C$ denotes the preconditions required to trigger a vulnerability, $A$ represents the action sequence that may activate it, $E$ specifies the expected buggy behavior, and $O$ defines the oracle assertions used to validate the bug.
Therefore, the goal of \consensustester is to construct an automated loop that follows the HDT paradigm and continuously discovers bugs by generating, instantiating, and validating vulnerability hypotheses.


\subsection{Consensus Protocols}
% %怎么写bug的级别？
%TLA+ 不写了吧
Consensus protocols mainly fall into two categories: Crash Fault-Tolerant (CFT) protocols, such as Paxos~\cite{Lamport1998ThePP,Lamport2001PaxosMS}, Fast Paxos~\cite{Lamport2006FastP}, EPaxos~\cite{Moraru2013ThereIM}, and Raft~\cite{Ongaro2014InSO}; and Byzantine Fault-Tolerant (BFT) protocols, such as PBFT~\cite{Castro1999PracticalBF}, HotStuff~\cite{Yin2018HotStuffBC}, DAG-Rider~\cite{Keidar2021AllYN}, Narwhal–Tusk~\cite{Danezis2021NarwhalAT}, and Bullshark~\cite{Giridharan2022BullsharkDB}.
CFT consensus protocols are primarily used in traditional banking systems and Internet companies for data transactions and cluster coordination, whereas BFT consensus protocols form the foundation of blockchain systems and on-chain transactions. These protocols are widely used in both research and production settings, with code repositories implemented in multiple programming languages, such as Go and Rust. The distinction is particularly important because BFT systems explicitly tolerate malicious behavior from participating nodes, whereas CFT systems do not; as a result, analyzing these two classes of protocols involves fundamentally different \textbf{Constraints}.

To analyze bugs in consensus protocols, it is essential to summarize \textbf{bug patterns} previously identified through extensive human expertise, in order to more effectively uncover bugs that violate system safety and correctness and may lead to erroneous transactions and data corruption (i.e., protocol-level \textbf{Logic Bugs}), rather than low-level \textbf{Implementation Bugs} such as out-of-memory errors, memory leaks, integer overflows, or simple logical flaws.
% 共识协议主要有Crash Fault-Tolerant (CFT),比如Paxos~\cite{Lamport1998ThePP,Lamport2001PaxosMS}, Fast Paxos~\cite{Lamport2006FastP}EPaxos ~\cite{Moraru2013ThereIM}Raft~\cite{Ongaro2014InSO},  Byzantine Fault-Tolerant (BFT) 比如PBFT~\cite{Castro1999PracticalBF},HotStuff~\cite{Yin2018HotStuffBC}, DAG-Rider~\cite{Keidar2021AllYN}, Narwhal-Tusk~\cite{Danezis2021NarwhalAT}, Bullshark~\cite{Giridharan2022BullsharkDB}。CFT 共识协议主要应用于传统的银行，互联网公司进行数据交易，集群调度； BFT共识协议主要是blockchain system和链上交易的基础。这些协议都广为使用的research 以及production level的代码仓库，使用多种不同的语言，比如Go Rust实现的。
% 想要分析共识协议的bug，
% 我们应该总结之前human knowledge查找的bug pattern,才能够更好地找到能够违背系统安全和正确性，导致交易数据错误的bug(Protocol-level logic bugs)
% 而不是 out-of-memory errors, memory leaks, integer overflows, or simple logical flaws （Low-level implementation bugs）.
% (这两者的区别是BFT系统允许有参与的节点恶意作恶，而CFT系统并不会考虑，因此分析这两类协议有不同的  constrains)。


% 分布式共识协议由 Crash Fault-Tolerant (CFT)~\cite{Lamport1998ThePP,Lamport2001PaxosMS,Moraru2013ThereIM,Ongaro2014InSO} and Byzantine Fault-Tolerant (BFT)~\cite{Castro1999PracticalBF,Lamport2011ByzantizingPB,Yin2018HotStuffBC,Danezis2021NarwhalAT} 组成。 

% 有非常多知名的代码仓库

% bug pattern很重要，

% 两者的区别是 constrains不同

%不同种类bug


%可解释性
% 多种不同的变成语言  go, Rust

%其中分布式协议需要满足的property主要是liveness和safety, 而诱发违背这两个property的原因主要分成五大类，Recovery & Execution Divergence
% Persistence & Monotonicity Violations
% Dependency & Topology Flaws
% Message Binding & Signature Violations
% Resource & Operational Visibility Violations

%凑字数的
The distributed protocols are primarily required to satisfy two fundamental properties: liveness and safety. Only the causes that lead to violations of these two properties constitute true protocol-level logic bugs. Such causes can be broadly categorized into five classes: (i) Recovery\&Execution Divergence, (ii) Persistence\&Monotonicity Violations, (iii) Dependency\&Topology Flaws, (iv) Message Binding\&Signature Violations, and (v) Resource\&Operational Visibility Violations.


\subsection{LLM-based Multi-Agent System}
The paradigm of LLM applications focuses on multi-agent architectures, which decompose complex tasks into collaborative, specialized roles. In this way, multi-agent systems enable more focused reasoning, mitigate hallucinations, and manage long-horizon objectives more effectively than single-agent models~\cite{Chen2023AgentVerseFM,Guo2024LargeLM,Yang2025AgentNetDE}. The efficacy of this paradigm also enables structured coordination mechanisms and communication protocols (see the formalization in Appendix~\ref{appendsec:multiagent}). 

Early dual-agent role-playing systems, such as CAMEL~\cite{Li2023CAMELCA} and ProAgent~\cite{Zhang2023ProAgentBP}, demonstrate that collaboration among agents can elicit step-by-step reasoning. More recent frameworks extend this idea by assigning explicit roles to multiple agents: MetaGPT~\cite{hong2023metagpt} and ChatDev~\cite{Qian2023ChatDevCA} employ manager, designer, and coder agents to simulate software engineering workflows; Magnetic-One~\cite{Fourney2024MagenticOneAG} and AG2~\cite{AG2} introduce a central coordinator that assigns tasks to different workers. OWL~\cite{Hu2025OWLOW} further decouples multi-agent systems, enabling effective transfer across domains to solve new tasks.
However, existing multi-agent systems still lack the ability to reason deeply about code patterns, especially to integrate domain knowledge of consensus protocol systems for efficiently discovering deep bugs.



\begin{algorithm}[!t]
\caption{\consensustester Workflow}
\label{alg:workflow}
\begin{algorithmic}[1]
    \STATE {\bfseries Input:} Repository Knowledge $\mathcal{R}$, Bug Patterns $\mathcal{P}$, Protocol Constraints $Con$
    \STATE {\bfseries Output:} Detected Bugs $\mathcal{B}$
    \STATE $\mathcal{S} \leftarrow \emptyset$ \COMMENT{Initialize global state}
    
    \WHILE{exploration budget remains}
        \STATE \textit{// --- Orchestrator Agent ---}
        \STATE $\mathcal{BE} \leftarrow \text{bug-exploitation}(\mathcal{B})$ \COMMENT{Analyze historical bugs found}
        \STATE $Z \leftarrow \text{state-analyzer}(\mathcal{S})$ \COMMENT{Summarize global state}
        
        \STATE \textit{// --- Strategy Agent ---}
        \STATE $AttS \leftarrow \text{Strategy}(\mathcal{BE}, Z, Con, \mathcal{P}, \mathcal{R})$ \COMMENT{Synthesize attack scenarios}
        % \STATE $\mathcal{S} \leftarrow \text{UpdateState}(S)$ \COMMENT{Sync scenarios to Orchestrator}
        \STATE \textbf{SendToOrchestrator}($AttS$) \COMMENT{Update $\mathcal{S}$}
        
        \STATE \textit{// --- TestGen Agent ---}
        \REPEAT
            \STATE $T \leftarrow \text{GenerateUnitTests}(AttS, \mathcal{R})$ \COMMENT{Repo knowledge-based gen unit test}
            \STATE $R \leftarrow \text{ExecuteAndAnalyze}(T)$
        \UNTIL{$R = \text{Success} \lor \text{MaxRetries}$}
        
        \IF{$R = \text{Success}$}
            %\STATE $\mathcal{B} \leftarrow \mathcal{B} \cup \{T\}$
            \STATE \textbf{SendToOrchestrator}($T$) \COMMENT{Update $\mathcal{S}$}
        \ENDIF
    \ENDWHILE
\end{algorithmic}
\end{algorithm}




\subsection{LLM-powered Bug Detection} 
%没写这个任务为何agent合适 intro第三段算是写了吧
% KernelGPT~\cite{Yang2025KernelGPTEK}分析了os kernel的bug
% \WhiteFox~\cite{Yang2023WhiteFoxWC} focus compiler
%这两个有domian knowledge
Recent advances in LLMs have significantly influenced automated bug detection and program analysis. Prior work has demonstrated that LLMs can assist developers in tasks such as bug localization, program repair, test generation, and vulnerability detection by leveraging their strong code understanding and natural language reasoning capabilities~\cite{feng2020codebert,chen2021evaluating,pearce2022asleep}. 


% LLMSAN~\cite{Wang2024SanitizingLL}, Scale~\cite{Wen2024SCALECS}, LLMDFA~\cite{Wang2024LLMDFAAD}, PrimeVul~\cite{Ding2025VulnerabilityDW} 都使用了LLM来帮助进行detect bugs.
% KNighter~\cite{Yang2025KNighterTS} 第一次使用agent系统来分析llm system code bugs.
% implementation bug
% However, existing LLM-powered d bug detection  techniques largely focus on localized code patterns (\textbf{implementation bugs}) and lack the ability to reason over complex, state-dependent behaviors that span multiple components or execution stages (\textbf{logic bugs}).  而且这些方法都无法利用multi-agent system符合HDT 范式来找到更加完备的bugs.
LLMSAN~\cite{Wang2024SanitizingLL}, SCALE~\cite{Wen2024SCALECS}, LLMDFA~\cite{Wang2024LLMDFAAD} and PrimeVul~\cite{Ding2025VulnerabilityDW} all leverage LLMs to assist in bug detection. KNighter~\cite{Yang2025KNighterTS} is the first to employ an agent-based system for analyzing bugs in LLM system code.
However, these existing LLM-powered bug detection techniques largely focus on localized code patterns (\textbf{implementation bugs}) and lack the ability to reason about complex, state-dependent behaviors that span multiple components or execution stages (\textbf{logic bugs}). Moreover, these approaches do not leverage multi-agent systems to combine with HDT paradigm to achieve more comprehensive bug discovery.



\section{Method}
\label{sec:method}


% 用了很久的库很多年的库etcd
%HDT也得写进来 % HDT align起来


\subsection{Motivation and Design Goal}
\label{sec:goal}
% 我们引入\consensustester,来融合consensus的domain knowledge with LLM-based mutiple agent system，来实现利用HDT更加完备detect Protocol-level logic bugs这一个根本性挑战. besides these, \consensustester 应该具有Interpretability，发现的bug具有详细的分析报告，方便用户debug;efficient,使用更少的tokens，发现的bug false postive少；effective，能够找到deep logic bugs; flexible,可以实现各个模块的解耦，方便即插即用推广到其他领域；complete，对CFT BFT系统都可以找到bug.

%方便用户使用 Interpretability
%token,false positive少  efficient
%focus on 深层次的bugs effective
%即插即用? flexible 解耦
%CFT BFT都行 complete
%HDT也得写进来 % HDT align起来

% 而且我们的方法应该尽量使用计算资源发现重要的logic bug，不能waste tokens on implementation bugs

We introduce \consensustester, which integrates domain knowledge of distributed consensus with an LLM-based multi-agent system to address the fundamental challenge of comprehensively detecting protocol-level logic bugs through HDT. Beyond this core capability, \consensustester is designed with several desirable properties. It is \textbf{{effective}}, capable of uncovering deep logic bugs; it provides  \textbf{{interpretability}}, producing detailed analysis reports for discovered bugs to facilitate debugging; it is \textbf{{efficient}}, requiring less communication and fewer tokens and yielding a low false-positive rate; it is  \textbf{{complete}}, applicable to both CFT and BFT systems; and it is \textbf{{flexible}}, featuring a modular, decoupled design that supports plug-and-play extension to other domains. 

Moreover, our method should prioritize using computational resources to discover important \textbf{logic bugs}, rather than wasting tokens on \textbf{implementation bugs}.



\subsection{Architecture}

% decoupling of workflow control, consensus scenarios properties, and test generation

%Orchestrator Agent bug-exploitation state-analyzer
%Strategy Agent Constraint Scenarios 
%TestGen Agent Loop

% 解耦 可以更好地即插即用  flexible
% interpretability %可解释性
%bug-exploitation能发现更深的bug effective
%Constraint complete
%实验系统证明了 efficient Sec~\ref{sec:exp}

% As shown in Figure~\ref{fig:overview}, \consensustester  comprises three core agents: (i) \bfagenta agent 负责workflow control，具有两个components, bug-exploitation 可以利用\consensustester 已经发现的bug更好地trigger出新的bug, state-analyzer可以维护另外两个agent的状态，并且防止\consensustester 总是尝试同样地bug场景， (ii) \bfagentb agent 具有两个componnet, Constraints-analyzer来分析bug attack Scenarios的限制,支持both CFT and BFT， 一个是attack Scenarios generator,来生成相应的场景来帮助发现深层次的bug；（iii） \bfagentc agent 使用一个reflection Loop 的componnet，自动生成，运行，分析，来得到降低false positive的unit tests，并且综合attack Scenarios生成用户友好的bug report 。

% 总之， 这种让三个agent，一个负责coordinate, 一个负责找到一个logic bug的场景，一个负责test generation,很好的解耦system, 实现了flexible和可拓展性。
% Sec~\ref{sec:exp} 证明了efficiency  and effectiveness of \consensustester.


As shown in Figure~\ref{fig:overview}, \consensustester is comprised of three core agents:
(i) \bfagenta agent is responsible for workflow control and consists of two components. The \texttt{bug-exploitation} component leverages bugs previously discovered by \consensustester to more effectively trigger new bugs, while the \texttt{state-analyzer} component maintains the states of the other two agents and prevents \consensustester from repeatedly exploring identical bug scenarios;
(ii) \bfagentb agent contains two components: a \texttt{constraints-analyzer} that analyzes the constraints of bug attack scenarios and supports both CFT and BFT protocols (\textbf{complete}), and an attack \texttt{scenario-generator} that produces protocol-specific scenarios to facilitate the discovery of deep logic bugs;
(iii) \bfagentc agent employs a \texttt{reflection-loop} component to automatically generate, execute, and analyze unit tests, reducing false positives and synthesizing user-friendly bug reports based on the generated attack scenarios  (\textbf{interpretability}).

Overall, this design, where one agent coordinates the workflow, one agent identifies protocol-level logic bug scenarios, and one agent focuses on test generation—effectively decouples system functionality, is composed of five components and enables \textbf{flexibility} and extensibility. Section~\ref{sec:exp} demonstrates the \textbf{efficiency} and \textbf{effectiveness} of \consensustester.


\subsection{Knowledge Library}
\label{sec:knowledgelibrary}
% Human Knowledge 爬虫

%哪里加入了human knowledge

% In bug patterns,我们通过爬取共识协议github repo上的确认是bug的issue， 并且进行了相应的总结，过滤掉了low-level implementation bugs, 因此我们引入了consensus的domain specific knowledge,来帮助\consensustester 理解protocol-level logic bugs 帮助生成更加attack 场景。 For constrains, 我们通过总结共识协议的setting和设计的不同，同时根据分析bug patterns 可以得到CFT和BFT系统trigger bug的不同：比如如果在cft系统中假定某个网络中的节点是malicious很容易违背系统的安全问题，但是实际上这并不practiical,而且会浪费计算资源。

% 通过 bug patterns和constrains, 我们为\consensustester 引入了需要的Knowledge Library。
% The details are shown in Appendix~\ref{appendsec:bugpattern} and Appendix~\ref{appendsec:constraints}.

For \texttt{bug patterns}, we crawl GitHub repositories of consensus protocols to collect issues that have been confirmed as bugs and summarize them accordingly, filtering out low-level implementation bugs. In this way, we incorporate domain-specific knowledge of consensus protocols to help \consensustester understand protocol-level logic bugs and generate more effective attack scenarios.
For \texttt{constraints}, we summarize the different settings and design choices of consensus protocols and, combined with the analysis of bug patterns, derive the distinct conditions under which bugs are triggered in CFT and BFT systems. For example, assuming malicious behavior of nodes in a CFT system can easily violate safety properties; however, such assumptions are impractical in real deployments and would unnecessarily waste computational resources.

By combining bug patterns and constraints, we introduce a knowledge library that provides the necessary domain knowledge for \consensustester. The details are presented in Appendix~\ref{appendsec:bugpattern} and Appendix~\ref{appendsec:constraints}.

\subsection{Memory}
\label{sec:memory}
%repo knowledge
% 多种不同的变成语言  go, Rust
%bug pattern   
% state很少
% 在我们的系统中，\consensustester需要维护的memory只有repo knowledge和bug patterns, bug patterns 已经在Sec~\ref{sec:knowledgelibrary}中详细描述了。而repo knowledge一般包含repo name, 共识协议的类型，协议的setting，实现的语言，测试的方法等，这样可以帮助我们的\consensustester 更加efficient支持各种不同编程语言和设计的协议。而且repo knowledge 和bug patterns被多个agnet需要用来做分析，并且非常小易于维护，因此我们将这两个放在memory中，在系统的环境中维护。

In our system, the memory maintained by \consensustester consists only of repository knowledge and bug patterns. Bug patterns are described in detail in Section~\ref{sec:knowledgelibrary}. Repository knowledge typically includes the repository name, the type of consensus protocol, protocol settings, implementation language, and testing methodologies. This information enables \consensustester to efficiently support consensus protocols implemented in different programming languages and with diverse designs.

Both repository knowledge and bug patterns are shared across multiple agents for analysis and are lightweight and easy to maintain (Appendix~\ref{appendsec:memory}). Therefore, we store them in the system memory and maintain them within the execution environment of \consensustester system.


\subsection{Workflow Details}
%sub-tasks  12  steps
%可解释性
%HDT也得写进来 % HDT align起来
%property + HDT
%无限badget

%step 1 bug exploitation subtask1
%prompt with 
%step 2 prompot state subtask2
%step 3发送 subtask3
% step 4 %读取constraints分析 subtask1
% prompot 整体
% step 5fuse 两个 subtask2
%step 6 attack场景 具体是如何扰动 subtask3
%step 7 发送 subtask4
%step 8 发动 subtask4 update
%step 9  prompt, generate test
% unit test
%step 10 prompt execute test
%step 11 analys 
%过 或者不过
%prompt可以合起来
%step 12 
% bug report 详细的有具体的信息
%发送 
%更新

% attack场景
% ：$C$: Preconditions（前置条件）—— 触发漏洞所需的初始状态  
% unit test
% ●$A$: Attack Actions（攻击动作）—— 触发漏洞的操作序列 
% execute and analyze 
% ●$E$: Expected Buggy Behavior（预期错误行为）—— 如果漏洞存在的可观察症状
% ●$O$: Oracle Assertions（断言）—— 验证漏洞的判定条件

% 在停止条件前或者消耗完badget前，As shown in Figure~\ref{fig:overview2}, the workflow of our \consensustester consists of 12 steps/subtasks.

% in step 1, \agenta agent进行其subtask1,使用prompt 进行bug exploitation，针对之前发现的bug继续分析，可以更好的guide \agentb基于之前的bugs得到新的deep bug.
% In step 2, \agenta agent进行其subtask2,使用prompt分析之前\agentb提出的attack场景和\agentc找到的bug 生成的global state.
% In step 3, \agenta agent 进行其subtask3， 将bug exploitation和state-analyzer后的结果发给\agentb,。

% In step4, \angetb agent进行其subtask1, 使用prompt读取constraints分析,
% in step 5,  \angetb agent进行其subtask2, fuse收到的 bug exploitation和state-analyzer后的结果
% in step 6, \angetb agent 进行其substaks3, 使用prompt利用step4和5的结果，生成attack scenarios,具体而言，可以根据具体的共识协议，模拟一个场景，控制node的加入，离线，carsh或者恢复，控制node发送message的顺序和冲突关系来生成容易产生bug的场景；
% in step 7, \agentb agent完成其subtask4,将生成的attack scenarios 发给\agenta agnet;
% in step 8, \agenta agent完成其subtask4,更新global state,将attack scenarios 发给\agentc agnet;
% in step 9, \agentc agent进行其subtask1根据收到的attack scenarios,利用prompt generate unit test.
% in step 10, \agentc agent进行其subtask2，使用prompt execute unit test;
% in step 11, \agentc agent进行其subtask3  利用prompt 分析生成的unit test和step10 的结果，如果成功，证明是找到consensus protocls的logc bugs,进入step 12,如果失败，则重新回到step 9 (循环的次数有一个上限 num\_iter, 当达到上限，证明这个attack scenario并不符合，需要重新生成)。
% in step 12, \agentc agent完成其substask4,将成功的bug 发给\agenta agent. 值得注意的事，虽然step9-11，看起来拆分的，其实可以使用一个prompt，让\agentc agent按照顺序完成三个subtask，实现循环。

Before reaching the termination condition or exhausting the budget, as shown in Figure~\ref{fig:overview2}, the workflow of \consensustester consists of 12 steps (subtasks).



In Step \circled{1}, \agenta performs its subtask O.1 and conducts bug exploitation using prompts. By analyzing previously discovered bugs, it provides better guidance for \agentb to derive new deep bugs based on prior bugs found from \agentc.
In Step \circled{2}, \agenta performs subtask O.2 and uses prompts to analyze attack scenarios proposed earlier by \agentb together with bugs discovered by \agentc, producing a global state.
In Step \circled{3}, \agenta performs subtask O.3 and sends the context information, i.e., the results of bug exploitation and state analysis, to \agentb.

In Step \circled{4}, \agentb performs subtask S.1 and analyzes protocol constraints using prompts.
In Step \circled{5}, \agentb performs subtask S.2 and employs the context from bug exploitation and state analysis.
In Step \circled{6}, \agentb performs subtask S.3 and, using prompts together with the results from Steps \circled{4} and \circled{5}, generates attack scenarios. Specifically, for a given consensus protocol, it simulates scenarios by controlling node behaviors such as joining, going offline, crashing, or recovering, as well as manipulating message sending orders and conflict relationships to produce bug-prone scenarios.
In Step \circled{7}, \agentb completes its subtask S.4 and sends the generated attack scenarios to \agenta.
In Step \circled{8}, \agenta completes its subtask O.4, updates the global state, and forwards the attack scenarios to \agentc.

In Step \circled{9}, \agentc performs its subtask T.1 and generates unit tests based on the received attack scenarios using prompts.
In Step \circled{10}, \agentc performs subtask T.2 and executes the generated unit tests using prompts.
In Step \circled{11}, \agentc performs subtask T.3 and analyzes the generated unit tests together with the execution results from Step \circled{10} using prompts. If successful, this indicates that a logic bug in the consensus protocol has been discovered, and the workflow proceeds to Step \circled{12}. Otherwise, the process returns to Step \circled{9}. The number of iterations is bounded by a maximum value, \texttt{num\_iter}; reaching this limit indicates that the current attack scenario is invalid and should be regenerated.
In Step \circled{12}, \agentc completes its subtask T.4 and reports the successfully discovered bug to \agenta. 
It is worth noting that although Steps \circled{9}–\circled{11} are described as separate subtasks, they can be implemented using a single prompt that enables the \agentc agent to sequentially complete the three subtasks in the reflection loop. %%songsa:这里对么？



%满足了property 和HDT
%每个都是四个subtasks

% 我们revisit HDT with our \consensustester, the attack scenario 符合 preconditions $C$, unit test generated by \agentc is the action sequence $A$,   the expected buggy
% behavior $E$ and the oracle assertions $O$ are all contained in step 10 and 11 of \agentc. 因此我们的\consensustester 很好的融合了consensus domain knowledge,以及HDT into our multi-agent systems.

% 总之， \agenta，\agentb, \agentc 都由4个substaks组成，prompt details of our agents could be found in Appnedix~\ref{appendsec:prompt}
% The workflow is also shown in Algorithm~\ref{alg:workflow}.

We revisit HDT with \consensustester. In our framework, the attack scenarios correspond to the preconditions $C$, the unit tests generated by \agentc constitute the action sequence $A$, and the expected buggy behavior $E$ and the oracle assertions $O$ are contained in Steps \circled{10} and \circled{11} of \agentc. As a result, \consensustester effectively integrates domain knowledge of consensus protocols and the HDT paradigm into a multi-agent system.
In summary, \agenta, \agentb, and \agentc are each composed of four subtasks. The prompt details of our agents are provided in Appendix~\ref{appendsec:prompt}. The overall workflow is illustrated in Algorithm~\ref{alg:workflow}. 



\begin{table*}[!th]
    \centering
    \caption{Overall results for the baselines (top of the table) and \consensustester (bottom of the table). The baselines are only able to identify low-level implementation bugs, whereas \consensustester can successfully uncover protocol-level logic bugs. In Total, the former correspond to implementation bugs, while the latter correspond to logic bugs. According to Sec~\ref{sec:goal}, fewer implementation bugs are preferred, whereas more logic bugs are desirable. The best results are highlighted in \textbf{bold}.}
    \vspace{0.5em}
    \label{tab:overall}
    \resizebox{\textwidth}{!}{
    \begin{tabular}{l|cccc|cccc|c}
        \toprule
        \textbf{Bug Type}
        & \multicolumn{4}{c}{Implementation Bug}
        & \multicolumn{4}{c}{\textbf{Logic Bug}}
        &  \\
        \cmidrule(lr){2-5} \cmidrule(lr){6-9} \cmidrule(lr){10-10}
         \textbf{Consensus Protocol} & \textbf{Raft} & \textbf{EPaxos} & \textbf{HotStuff} & \textbf{BullShark}
        & \textbf{Raft} & \textbf{EPaxos} & \textbf{HotStuff} & \textbf{BullShark}
        & \textbf{Total} \\
        \midrule
        % \multicolumn{10}{l}{\textbf{Baselines}} \\
        \midrule
        GPT-5.2~\cite{openai2025gpt52} & 4 & 7 & 11 & 0 & 0 & 0 & 0 & 0 & 22,0 \\
        Gemini~3.0~Pro~Preview~\cite{google2025gemini3} & 0 & 0 & 0 & 0 & 0 & 0 & 0 & 0 & 0,0 \\
        Claude~Sonnet~4.5~\cite{anthropic2025claude45} & 0 & 0 & 0 & 0 & 0 & 0 & 0 & 0 & 0,0 \\
        Qwen3~Coder~480B~A35B~\cite{qwen2025qwen3} & 0 & 3 & 0 & 0 & 0 & 0 & 0 & 0 & 3,0 \\
        \midrule
        \rowcolor{gray!15} \textbf{Total (Baselines)} & 4 & 7 & 11 & 0 & 0 & 0 & 0 & 0 & 22,0 \\
        \midrule
        % \multicolumn{10}{l}{\textbf{\consensustester}} \\https://api.semanticscholar.org/CorpusID:273877854
        \midrule
        GPT-5.2~\cite{openai2025gpt52} & 0 & 0 & 0 & 0 & 1 & 3 & 3 & 1 & 0,8 \\
        Gemini~3.0~Pro~Preview~\cite{google2025gemini3} & 0 & 0 & 0 & 0 & 1 & 6 & 3 & 1 & \textbf{0,11} \\
        Claude~Sonnet~4.5~\cite{anthropic2025claude45} & 0 & 0 & 0 & 0 & 0 & 4 & 1 & 1 & 0,6 \\
        Qwen3~Coder~480B~A35B~\cite{qwen2025qwen3} & 0 & 0 & 0 & 0 & 1 & 6 & 2 & 0 & 0,9 \\
        \midrule
         \rowcolor{blue!10} \textbf{Total (\consensustester)} & 0 & 0 & 0 & 0 & 1 & 9 & 4 & 1 & \textbf{0,15} \\
        \bottomrule
    \end{tabular}
    }
\end{table*}


\subsection{Communication Mechanism} 
%knowledge +降低了消耗，毕竟要记录coor
%global state
% MCP (Model Context Protocol) memory调用其实也很少
% state很少

% Communication in \consensustester contains memory and global state and 共识协议的代码仓。
% 如Sec~\ref{sec:memory}提到的，memory很小，而且在整个运行区间agent获取一次即.
% 我们通过利用\agenta上的global state,这样可以让\agentb 无需担心生成多轮attack scenarios和bugs，可以轻松的通过\agenta获取状态，\agentc则可以通过\agenta获得attack scenarios。因此\agentb 和\agentc不需要考虑上下文信息，而且可以专注于本身任务,而\agenta保持整个系统中的状态，有效管理长期的目标。 state本身也很succinct, 而且其传输只是\agenta和\agentb, \agentc之间的，不涉及\agentb \agentc之间的沟通，而这种传输是维持系统状态必须，无法减少。
% in \consensustester, agent可以通过Model Context Protocol （MCP）获取共识协议代码仓。

% 因此我们的设计，让整个系统非常的简洁，通信只有必要的信息，commucnation overhead很小。


Communication in \consensustester involves the memory, the global state, and the consensus protocol code repository.

As discussed in Section~\ref{sec:memory}, the memory is lightweight and is accessed by the agents only once during the entire execution. By maintaining the global state within \agenta, \agentb does not need to track previously generated attack scenarios or discovered bugs and can easily retrieve the required state information from \agenta. Similarly, \agentc obtains the attack scenarios through \agenta. As a result, \agentb and \agentc do not need to manage contextual information and can focus solely on their respective tasks, while \agenta maintains the system-wide state and effectively manages long-horizon objectives.
The global state itself is succinct, and its transmission occurs only between \agenta and \agentb or \agentc, without direct communication between \agentb and \agentc. Such communication is necessary to maintain system consistency and cannot be further reduced. In \consensustester, agents can access the consensus protocol code repositories via the Model Context Protocol (MCP)~\cite{anthropic2024mcp}.

Overall, this design keeps the system concise, with communication limited to essential information, resulting in low communication overhead.
%songsa: state本身也很succinct, 有大小不;很小很小




\begin{table}[!t]
    \centering
    \caption{Results for bug detection in ablation studies. We compare the original \consensustester with variants in which the corresponding components are removed.}
    \vspace{0.5em}
    \label{tab:ablation_study}
\resizebox{\linewidth}{!}{
    \begin{tabular}{lccccc}
        \toprule
        \textbf{Methods}
        & \textbf{Raft} & \textbf{EPaxos} & \textbf{HotStuff} & \textbf{BullShark}
        & \textbf{Total} \\
        \midrule
        \consensustester & 1 & 6 & 3 & 1 & 11 \\
        w/o bug-exploitation & 0 & 0 & 3 & 0 & 3 \\
        w/o state-analyzer & 0 & 0 & 0 & 0 & 0 \\
        w/o constraints-analyzer & 0 & 0 & 1 & 0 & 1 \\
        w/o scenario-generator & 0 & 0 & 0 & 0 & 0 \\
        w/o reflection-loop & 0 & 0 & 0 & 0 & 0 \\
        \bottomrule
    \end{tabular}
    }
\end{table}



\section{Experiments}
\label{sec:exp}

\subsection{Experiment Setup} 



%解释下目前没有别的baselines了
% Claude Sonnet 4.5~\cite{anthropic2025claude45}
% Gemini 3.0 Pro Preview~\cite{google2025gemini3}
% GPT-5.2~\cite{openai2025gpt52}
% %open-source
% Qwen3 Coder 480B A35B~\cite{qwen2025qwen3}

% 由于目前的利用llm进行找bug的论文要么缺少和multi agent system的结合，要么完全不focus on 深层次的共识协议的logic bug。因此我们我们修改ReAct~\cite{yao2022react} with 读取代码的tools, implement it with current先进的code 大语言模型作为baselines，我们使用了三个商业大语言模型：Claude Sonnet 4.5~\cite{anthropic2025claude45}
% Gemini 3.0 Pro Preview~\cite{google2025gemini3}
% GPT-5.2~\cite{openai2025gpt52}
% 和一个开源的大语言模型 Qwen3 Coder 480B A35B~\cite{qwen2025qwen3}。
\textbf{Baselines.}
Existing studies that leverage LLMs for bug detection either lack integration with multi-agent systems or do not focus on deep logic bugs in consensus protocols. To address these limitations, we adapt ReAct~\cite{yao2022react} by augmenting it with code-reading tools, and implement it using state-of-the-art code-oriented large language models as baselines.
Specifically, we evaluate three commercial LLMs, Claude Sonnet 4.5~\cite{anthropic2025claude45}, Gemini 3.0 Pro Preview~\cite{google2025gemini3}, and GPT-5.2~\cite{openai2025gpt52}, as well as one open-source model, Qwen3 Coder 480B A35B~\cite{qwen2025qwen3}.


% codebase: Append~\ref{appendsec:codebase}%appendix: 代码仓star很多  % 多种不同的变成语言  go, Rust %reseach 还是production
%多少个iteration testgeg 50
%3个小时，iteration? overall
%temperature 0.2

% 所有大语言模型的temperature我们设定为0.2, 我们每次实验都设定\consensustester和baseline运行三小时, 对于\consensustester而言reflection loop的  num\_iter为50.我们针对2个BFT协议和2个CFT协议的的protocol的四个popular code repo进行了实验, Raft (etcd~\cite{etcdioraft}, EPaxos (efficientEPaxos~\cite{efficientepaxos})
% HotStuff (relabHotStuff~\cite{suimystenlabs}), BullShark~\cite{suimystenlabs}
% 其中etcd和sui是prodcution 的代码仓

\begin{figure}[!t]
  \centering
  \includegraphics[width=0.8\linewidth]{tables/bug_class_distribution.png}
  \caption{Distribution of discovered protocol-level logic bugs of \consensustester by category.}
  \label{fig:bugclass}
\end{figure}



\textbf{Implementation Details.}
For our implementation, we access all models via APIs, eliminating the need for GPUs. For all large language models, we set the temperature to 0.2. In each experiment, both \consensustester and the baselines are allowed to run for three hours. For \consensustester, the number of reflection loop iterations (\texttt{num\_iter}) is set to 50. The implementation and artifact details are in Appendix~\ref{appendsec:artifact}.


\textbf{Code Repositories.} We evaluate our approach on four popular protocol code repositories, covering two BFT protocols and two CFT protocols, including Raft (etcd~\cite{etcdioraft}), EPaxos (efficientEPaxos~\cite{efficientepaxos}), HotStuff (relabHotStuff~\cite{suimystenlabs}), and BullShark (sui~\cite{suimystenlabs}). Among these, etcd and sui are production-level code repositories. The details of these repositories are shown in Appendix~\ref{appendsec:codebase}.



\begin{figure*}[!t]
    % \captionsetup{skip=2pt}
    \centering
    \subfloat{
        \includegraphics[width=0.45\linewidth]{tables/total_tokens.png}
        \label{fig:totaltokens}
    }
     \hspace{0.04\linewidth}
    % \vspace{-1em}
    \subfloat{
    \includegraphics[width=0.45\linewidth]{tables/tokens_per_bug.png}
        \label{fig:ablationtokens}
    }
    \caption{The left (a) shows the total token consumption in the ablation studies.
The right (b) shows the average token consumption for each discovered bug in the ablation studies.}
\end{figure*}



\subsection{Overall Results}
\label{sec:expoverall}

% false positive 20-30% 说理性的
%HDT导致我们xx
%开源框架和闭的对比 类似owl写很多
%gemini3.0最好
% 强调bug level
% 发现很多不同的bug
%有重复的
% 两种bug
%具体分析
% false positive 

% Table~\ref{tab:overall} shows the overall results of our experiments from \consensustester against baselines. 从结果中,我们可以看到, baselines不管使用什么llm都只能发现简单的implementaion的bug, 而我们的\consensustester 的设计可以在任何llm作为基底的情况下都能够发现protocol的logic bugs,而不会发现 implementaion bug.  由于implementaion levle的bug实际上可以通过静态分析发现,只会导致integer overflow 或者内存溢出等问题, logic bug更加隐秘可能导致系统数据和交易的安全问题. 因此我们\consensustester focus on deep logic bugs. The details of the two bug types are in Appendix~\ref{appendsec:bugexample}. ,实验结果证明了我们的\consensustester可以使我们设计的multiagent system 更好focus我们需要的detect deep bug的需求而非implementation bug,说明了我们设计的有效性.

% 具体而言,我们发现即便融合了ReAct和llm,Gemini 3.0 Pro Preview和Claude Sonnet 4.5依然无法发现任何implementation bug,GPT-5.2  在baseline中最为有效可以在三个repo上发现22个bug,而Qwen3 Coder 480B A35B 可能由于其对code 功能的优化也可以在Epaxos的代码库上发现3个bug(这三个bug也可以被GPT-5.2发现).因此,in total, 我们实现的baseline一共发现了22个zero-day bug,这也是我们本文对共识protocol implementation修复的贡献. 


% Compared to the results of baseline on Gemini 3.0 Pro Preview (0 bugs),  我们的\consensustester 在上使用Gemini 3.0 Pro Preview效果最好,可以在包含广泛使用的production-level的四个repo上发现11个bug,这证明了我们的方法的有效性.\consensustester在其余三个llm都能发现logic bug的结果,证明了我们方法的有效性.由于使用四种llm会发现部分相同的bug,我们一共可以发现15个zero-day logic bug,可以大大提高分布式领域的安全性.这也证明了我们框架和HDT结合的有效性.

% besides these,我们还分析了\consensustester中\agentc agent使用四个llm给出的所有的46,bug report,其中一个有34个是真实的logic bug,我们的\consensustester仅有26.1\%的false positive.这证明了使用我们的\consensustester 找到的bug并不需要浪费大量地人工去filter过度的bug report,非常的高效.


Table~\ref{tab:overall} presents the overall experimental results of \consensustester compared with the baselines. From the results, we observe that regardless of the underlying LLM, the baselines are only able to identify simple implementation-level bugs, whereas our \consensustester, \textbf{by design}, consistently discovers protocol-level logic bugs and does not report implementation bugs.

Implementation-level bugs can typically be detected through static analysis and often result in issues such as integer overflows or memory violations. In contrast, logic bugs are significantly more subtle and may compromise system correctness, data consistency, or transaction security, leading to severe real-world consequences. Therefore, \consensustester explicitly focuses on deep protocol logic bugs rather than implementation-level issues. Detailed examples of the two bug types are provided in Appendix~\ref{appendsec:bugexample}. The experimental results demonstrate that our multi-agent design enables \consensustester to effectively focus on the detection of deep logic bugs, saving the computational resources and validating the effectiveness of our approach.



More specifically, we find that even when augmented with ReAct design, baselines using Gemini 3.0 Pro Preview and Claude Sonnet 4.5 fail to identify any implementation-level bugs. GPT-5.2 is the most effective baseline model, detecting 22 bugs across three repositories. Qwen3 Coder 480B A35B, possibly benefiting from its optimization for code understanding, is able to detect three bugs in the EPaxos repository, all of which are also identified by GPT-5.2. In total, the baselines uncover 22 zero-day implementation bugs, which we report as an additional contribution toward improving robustness of existing consensus implementations.

In contrast, \consensustester achieves the best performance when built upon Gemini 3.0 Pro Preview, identifying 11 protocol-level logic bugs across four widely used repositories, including production-level systems. This result is particularly notable given that the Gemini-based baseline fails to detect any bugs. Moreover, \consensustester consistently discovers logic bugs when instantiated with each of the other three LLMs, further demonstrating the robustness and generality of our approach. As some bugs are discovered by multiple LLM-backed instances, \consensustester identifies a total of 15 unique zero-day logic bugs, substantially enhancing the security of distributed consensus systems. These results further validate the effectiveness of integrating our framework with the HDT paradigm.

\textbf{Practical Usage.} We further analyze all the 46 bug reports generated by the \agentc agent in \consensustester across the four LLMs. Among them, 34 reports correspond to real protocol-level logic bugs, resulting in a false positive rate of \textbf{only 26.1\%}. This low false positive rate indicates that \consensustester does not require extensive manual filtering, making it a highly efficient and practical tool for engineers and researchers to discover deep logic bugs in real-world systems.




% \begin{figure}[!t]
%   \centering
%   \includegraphics[width=0.88\linewidth]{tables/total_tokens.png}
%   \caption{Total token consumption in the ablation studies.}
%   \label{fig:totaltokens}
% \end{figure}


% \begin{figure}[!t]
%   \centering
%   \includegraphics[width=0.88\linewidth]{tables/tokens_per_bug.png}
%   \caption{Average token consumption for each discovered bug in the ablation studies.}
%   \label{fig:ablationtokens}
% \end{figure}








\textbf{Bug Summary.}
Based on the categories of consensus logic bugs, we further analyze the 15 bugs discovered by \consensustester. As shown in Figure~\ref{fig:bugclass}, violations are observed across all bug categories, demonstrating that \consensustester is capable of triggering a wide range of deep protocol-level bugs.
Specifically, Recovery \& Execution Divergence bugs account for four cases (26.7\%), tying with Persistence \& Monotonicity Violations as the most prevalent category. In addition, Message Binding \& Signature Violations, Resource \& Operational Visibility Violations, and Dependency \& Topology Flaws account for 2, 2, and 3 bugs, respectively. Detailed examples of logic bugs from different categories are in Appendix~\ref{appendsec:bugclassexample}.

% 基于共识logic bug的类别,我们进一步分析了我们发现的15个bugs. 如图~\ref{fig:bugclass} 所示,我们发现我们可以在所有的类别上都发现bug的violation,这证明了\consensustester可以trigger出各种深层次的protocol-level bugs. 其中Recovery\&Execution Divergence bug有4个,占比26.7\%和Persistence\&Monotonicity Violations 并列最多; Message Binding\&Signature Violations, Resource\&Operational Visibility Violations和Dependency\&Topology Flaws,我们分别发现了2,2,3个bug. The details of 不同类别的logic bug are shown in Appendix~\ref{appendsec:bugclassexample}.


%需要一个大饼图
% Recovery\&Execution Divergence, Persistence\&Monotonicity Violations, Dependency\&Topology Flaws, Message Binding\&Signature Violations, and Resource\&Operational Visibility Violations
% Appendix~\ref{appendsec:bugclassexample}.



% \begin{figure}[!th]
%   \centering
%   \includegraphics[width=0.9\linewidth]{tables/bugs_per_hour.png}
%   \caption{Ablation hour}%
%   %结果是尖括号
%   \label{fig:ablationhour}
% \end{figure}




\subsection{Ablation Studies}
% bug-exploitation 
%  state-analyzer
%  constraints-analyzer
%  scenario-generator
%  reflection-loop


%强调去掉了一些components，我们的也能发现low-level bugs！！！！对比baselines

 % Gemini 3.0 Pro Preview~\cite{google2025gemini3}因为最好
%five components




% 在这个subsection, 我们进行\consensustester 五个componnet,bug-exploitation , state-analyzer,constraints-analyzer,和  reflection-loop 对发现bug和发现bug所需要计算资源的影响的分析,以及三个agent功能的测试.following之前同样地设置, 我们进行了ablation studies, 并且选取Gemini 3.0 Pro Preview作为我们的llm model由于其在Sec~\ref{sec:expoverall}中取得最好的效果.

% Table~\ref{tab:ablation_study}总结了abaltion study中找的的bug的results. 实验结果证明只有\consensustester才能发现最多的bug,而如何任何一个component都会导致实验效果的下降.而Figure~\ref{fig:ablationtokens} 则统计了在ablation study的情况下,发现每个bug所需要的token数量.

% 具体而言,我们发现当去掉bug-exploitation模块后,虽然方法依然能够在hotstuff上发现一些bug, 但是无法在其他的repo上发现logic bugs,而且无法利用之前的bug results进一步发现深入的bug,浪费了大量地tokens. The details of bugs triggered by bug-exploitation are shown in Appendix~\ref{appendsec:deepbug}. 这证明了我们的方法利用bug-exploitation,可以更好地发现更多的deep bugs.

% 而当我们去除 state-analyzer 模块,其无法发现任何bug,我们进行深入的研究发现,\consensustester在这种情况下失去上下文的context,不断重试场景,耗费tokens,无法发现logic bugs.

% 如果remove constraints-analyzer,虽然在BFT系统中HotStuff上发现一个bug,但是会导致在针对CFT系统的测试中生成false positives, 在我们的实验中,一共生成了107个CFT的bug reports,经过分析不是真的logic bugs.因此融入constraints-analyzer into \consensustester是必要的. Differences between BFT and CFT bugs are in Appendix~\ref{appendsec:bftbug} and Appendix~\ref{appendsec:cftbug}.

% 去除scenario-generato,我们的方法将会回退到baselines,如Appendix~\ref{appendsec:totaltoken} 所示,\consensustester在这种场景下会消耗最多的tokens,发现众多的implementation level的bugs,却无法发现任何logic bugs. 我们在这组实验中一共发现了29zero-day implementation levle bugs,证明了我们系统其他模块在自动迭代系统发现bug的有效性.

% 而去除reflection-loop 将会导致\consensustester无法生成正确的unit test文件,因而导致无法迭代生成logic bugs而浪费tokens,降低efficiency.


% 总之,我们的\agenta agent可以使用bug-exploitation和state-analyzer模块节约context和communcation沟通,\agentb agent可以生成发现deep bug的有效attack场景,而且使用 constraints-analyzer 更好应对不同类别protocols,\agentc agent 可以利用reflection-loop 生成正确用例,因此\consensustester保证了effective, interpretability, efficient,  complete,flexible .

% 58604702 43768170 79622940         41603,761               233076979.5             46911748

In this subsection, we analyze the impact of the five components of \consensustester on both bug discovery effectiveness and the computational resources required to discover bugs. Following the same experimental settings as before, we conduct ablation studies and select Gemini 3.0 Pro Preview as the underlying LLM, as it achieves the best overall performance under overall experiments in Section~\ref{sec:expoverall}.

\textbf{Summary.} Table~\ref{tab:ablation_study} summarizes the bug discovery results in the ablation study. The results show that only the complete \consensustester is able to discover the largest number of bugs (11 bugs), while removing any single component leads to degraded performance (at most 3 bugs). 
Figure~\ref{fig:totaltokens} shows the total token consumption in our ablation studies. Although removing the \texttt{bug-exploitation} and \texttt{reflection-loop} components can slightly reduce token usage (43.8M and 46.9M tokens, respectively, compared to 58.6M for the original configuration), this significantly hampers the discovery of useful bugs. Removing any of the other components results in higher token consumption, with the most significant increase observed when removing the \texttt{scenario-generator} (we will discuss later). Combined with Table~\ref{tab:ablation_study}, these results demonstrate that our modular design enables \consensustester to discover more meaningful bugs with fewer tokens, validating the effectiveness of our proposed components. Figure~\ref{fig:ablationtokens} further reports the number of tokens required to discover each bug under different ablation settings.

\textbf{bug-exploitation}. Removing \texttt{bug-exploitation} component is still able to discover a small number of bugs (3 bugs) in HotStuff, but fails to identify logic bugs in other repositories. Moreover, it cannot further explore deeper bugs based on previously discovered ones, resulting in significant token waste. Detailed examples of deeper bugs triggered through bug exploitation are provided in Appendix~\ref{appendsec:deepbug}. These results demonstrate that bug-exploitation enables \consensustester to more effectively discover deeper protocol-level bugs.

When the \texttt{state-analyzer} component is removed, \consensustester fails to discover any bugs. Further analysis shows that without this component, the system loses essential execution context, repeatedly retries similar scenarios, consumes excessive tokens, and is unable to uncover logic bugs.

Removing the \texttt{constraints-analyzer} component allows the system to discover a single bug in HotStuff (BFT protocol), but leads to a large number of false positives when testing CFT systems. In our experiments, this configuration generates 107 CFT bug reports, none of which correspond to real logic bugs after manual inspection. This highlights the necessity of integrating the constraints-analyzer into \consensustester. Detailed differences between BFT and CFT bugs are discussed in Appendix~\ref{appendsec:bftbug} and Appendix~\ref{appendsec:cftbug}.

When the \texttt{scenario-generator} component is removed, \consensustester regresses to baseline behavior. As shown in Figure~\ref{fig:totaltokens}, this configuration incurs the highest token consumption (233.1M tokens, 3.98 times of the original), discovers numerous implementation-level bugs, yet fails to identify any protocol-level logic bugs. In this setting, we uncover 29 zero-day implementation-level bugs, further demonstrating the effectiveness of the remaining components in enabling iterative bug discovery.


Finally, removing the \texttt{reflection-loop} component prevents \consensustester from generating correct unit test files, thereby breaking the iterative discovery process for logic bugs. This results in wasted tokens and reduced efficiency.

In summary, \agenta agent leverages the bug-exploitation and state-analyzer component to reduce context usage and communication overhead; \agentb agent generates effective attack scenarios for discovering deep bugs and adapts to different protocol types through the constraints-analyzer; and \agentc agent relies on the reflection-loop component to generate correct test cases. Together, these components ensure that \consensustester achieves the key properties of interpretability, efficiency, completeness, and flexibility.

%具体数字



% \label{appendsec:totaltoken}

%去掉了xx和xx一样

% Or 功能：节约context, communication的沟通
% St : 功能：生成靠谱的attack场景
% TestGen:功能：生成正确的用例，运行测试 unit test



% Test-time Scaling.  reflection loop里面



% effective, interpretability, efficient,  complete,flexible 关键!









\section{Limitations and Discussions}
\label{sec:limitation}
% limitation discussion
%SFT DPO等等
 % 解耦 可以更好地即插即用
 % transfer domain knowledge 没写
  % future work

 % 在本文中，我们针对consensus protocol的bug 提出了一个multiple agent system \consensustester,虽然我们效果已经很显著了，发现了很多zero-day logic bugs for consensus. 但是我们的方法依然是需要一些human knowledge，未来应该继续提高其自动化程度。

 % 由于共识的bug很难查找，因此我们目前只有15条数据，但是未来我们会extend我们的bug exploitation,并且将我们的实验部署到更多的consensus protocol代码仓上，收集足够的数据，利用Supervised Fine-Tuning（SFT）~\cite{Wei2021FinetunedLM,Ouyang2022TrainingLM}提高我们的系统，而且可以基于此生成 trajectory 利用 Direct Preference Optimization  (DPO)~\cite{Rafailov2023DirectPO}来微调我们到的agent system.

 % 而且未来我们将会推广\consensustester 利用HDT和我们解耦的设计
 % 到protocol和实现都很复杂的debug领域，比如并发控制和smart contract，都容易引发现实中严重经济安全问题。

% In this paper, we propose \consensustester, a multi-agent system for discovering bugs in consensus protocols.
Although our system has demonstrated significant effectiveness and has uncovered many zero-day logic bugs in consensus implementations, it still requires a certain amount of human knowledge. In the future, we aim to further improve its level of automation.
Due to the inherent difficulty of identifying consensus bugs, we find 15 bug instances. In future work, we plan to extend our bug exploitation techniques and deploy our experiments on a broader range of consensus protocol codebases to collect sufficient data. With more data, we intend to apply Supervised Fine-Tuning~\cite{Wei2021FinetunedLM,Ouyang2022TrainingLM} to enhance our system, and further generate trajectories to finetune our agent system using Direct Preference Optimization~\cite{Rafailov2023DirectPO}.

Moreover, we plan to generalize \consensustester, leveraging HDT and our decoupled design, to other debugging domains where both protocols and implementations are highly coupled and complex, such as concurrency control~\cite{Tan2020CobraMT,Clark2024ValidatingDS} and smart contracts~\cite{Sun2024GPTScanDL}, which are prone to causing severe economic and security issues in real-world systems.


\section{Conclusion}
\label{sec:conclusion}
% HDT
% effective, interpretability, efficient,  complete,flexible
% 15个bug
% 多种不同的变成语言  go, Rust %reseach 还是production

% 在本文中，我们将specific domain knowledge和llm-based multiple agent system进行了结合，利用HDT范式，提出了\consensustester 来detect consensus protocol-level logic bug。我们的方法针对不管是research还是production level的code repo都可以通过生成特定的攻击场景，发现相应的深层次的对实际高威胁的bug，可以支持多种编程语言，大大提高了分布式系统领域的安全性，并且明显优于baselines。我们方便解耦的设计，不仅仅满足了 effective, interpretability, efficient,  complete,flexible等property，而且还可以未来推广拓展到其他需要domain knowledge的bug detections.

In this paper, we combine specific domain knowledge with an LLM-based multi-agent system and, following the HDT paradigm, propose the \consensustester framework to detect consensus protocol–level logic bugs. Our approach is applicable to both research prototypes and production-level code repositories. By generating targeted attack scenarios, \consensustester is able to uncover deep, high-impact vulnerabilities that pose serious real-world threats. The system supports multiple programming languages, significantly enhancing the security of distributed systems, and demonstrates clear advantages over existing baselines.
Benefiting from our cleanly decoupled design, our multi-agent framework not only satisfies key properties such as effectiveness, interpretability, efficiency, completeness, and flexibility, but also provides a strong foundation for future extension to other domain knowledge–intensive bug detection tasks.




%related works再来点
% 过一次大模型写作

% figure 1似乎没强调implementation bug

% table如何改 强调total
%单纯大模型，我们专门找deep bug设计的

%加粗 实验里的文字,实验里的数字不够
%可以写更多,压缩空间 vspace就行
%ablation修改 更多内容，不然显得不是很有效
 % efficiency and effectiveness  没写 in ablation
 



% \section*{Acknowledgements}
\clearpage

\section*{Impact Statement}
This paper presents work whose goal is to advance the field of 
Machine Learning. There may be some potential societal consequences 
of our work, none which we feel must be specifically highlighted here.

% In the unusual situation where you want a paper to appear in the
% references without citing it in the main text, use \nocite
% \nocite{langley00}

\bibliography{agora}
\bibliographystyle{icml2025}


%%%%%%%%%%%%%%%%%%%%%%%%%%%%%%%%%%%%%%%%%%%%%%%%%%%%%%%%%%%%%%%%%%%%%%%%%%%%%%%
%%%%%%%%%%%%%%%%%%%%%%%%%%%%%%%%%%%%%%%%%%%%%%%%%%%%%%%%%%%%%%%%%%%%%%%%%%%%%%%
% APPENDIX
%%%%%%%%%%%%%%%%%%%%%%%%%%%%%%%%%%%%%%%%%%%%%%%%%%%%%%%%%%%%%%%%%%%%%%%%%%%%%%%
%%%%%%%%%%%%%%%%%%%%%%%%%%%%%%%%%%%%%%%%%%%%%%%%%%%%%%%%%%%%%%%%%%%%%%%%%%%%%%%
\newpage
\appendix
\onecolumn

\begin{center}
    {\LARGE \bf Appendix}
    % \vspace{0.5cm}
\end{center}
\begin{quote}
    \hypersetup{linkcolor=black} % 临时将链接设为黑色，看起来更像正式出版物
    \textbf{Contents}
    \begin{itemize}
    \item \hyperref[appendsec:artifact]{Section \ref{appendsec:artifact}: Artifact Details}
        \item \hyperref[appendsec:multiagent]{Section \ref{appendsec:multiagent}: Mathematical Modeling of the Multi-Agent System}
        \item \hyperref[appendsec:bugpattern]{Section \ref{appendsec:bugpattern}: Bug Pattern Details}
        \item \hyperref[appendsec:constraints]{Section \ref{appendsec:constraints}: Constraints Details}
        \item \hyperref[appendsec:memory]{Section \ref{appendsec:memory}: Memory Details}
        \item \hyperref[appendsec:codebase]{Section \ref{appendsec:codebase}: Codebase Details}
        % \item \hyperref[appendsec:totaltoken]{Section \ref{appendsec:totaltoken}: Ablation Stufy for Total Token Usage}
        \item \hyperref[appendsec:bugexample]{Section \ref{appendsec:bugexample}: Two Types of Bug Example}
        % \begin{itemize}
        %     \item \hyperref[appendsec:implementationbug]{F.1 Implementation Bug}
        %     \item \hyperref[appendsec:logicbug]{F.2 Logic Bug}
        % \end{itemize}
        \item \hyperref[appendsec:bugclassexample]{Section \ref{appendsec:bugclassexample}: Bug Example for Five Classes}
        \item \hyperref[appendsec:prompt]{Section \ref{appendsec:prompt}: Prompt Example}
    \end{itemize}
\end{quote}
% \vspace{1cm}
\hrule
% \vspace{1cm}
% --- 手动附录目录结束 ---
%%%%%%%%%%%%%%%%%%%%%%%%%%%%%%%%%%%%%%%%%%%%%%%%%%%%%%%%%%%%%%%%%%%%%%%%%%%%%%%
%%%%%%%%%%%%%%%%%%%%%%%%%%%%%%%%%%%%%%%%%%%%%%%%%%%%%%%%%%%%%%%%%%%%%%%%%%%%%%%

\section{Artifact Details}
\label{appendsec:artifact}
We have uploaded the codebase for our paper, which is available in the supplementary materials. Setting up the environment is straightforward with the provided README file. In the README, you will find all the necessary guidelines to use the bash scripts to reproduce the main results from our experiments.

\section{Mathematical Modeling of the Multi-Agent System}
\label{appendsec:multiagent}
In a formal sense, we characterize an autonomous LLM-powered entity as a quintuple $\mathcal{E} = \langle \mathcal{B}, \mathcal{P}, \mathcal{A}, \Omega, \Gamma \rangle$. Here, $\mathcal{B}$ represents the core linguistic brain (the LLM), $\mathcal{P}$ signifies the perceptual manifold of the target environment, and $\mathcal{A}$ denotes the repertoire of executable operations (tools). The mapping function $\Omega: \mathcal{P} \to \mathcal{A}$ acts as the decision logic, while $\Gamma$ encapsulates the persistent memory state. The operational trajectory of such an entity follows a recursive observe-deliberate-act cycle: at any discrete temporal index $k$, the entity ingests a sensory input $p_k \in \mathcal{P}$, leverages $\mathcal{B}$ to synthesize a strategic intent based on $\Omega$, and dispatches an intervention $a_k \in \mathcal{A}$ to modulate the system state. To tackle the intricate state spaces of consensus implementations, we extend this individual paradigm into a collaborative collective $\mathcal{G} = \{ \mathcal{E}_1, \mathcal{E}_2, \dots, \mathcal{E}_m, \mathcal{X} \}$. In this formulation, each $\mathcal{E}_i$ represents a specialized functional unit, and $\mathcal{X}$ defines the coordination medium that facilitates information exchange and role synchronization. This multi-agent architecture transcends the cognitive boundaries of isolated models by decomposing complex protocol verification into modular, inter-dependent reasoning sub-tasks, a methodology that has proven more robust than monolithic approaches in recent literature, e.g., CAMEL~\cite{Li2023CAMELCA}, MetaGPT~\cite{hong2023metagpt}, OWL~\cite{Hu2025OWLOW}.


\section{Bug Pattern Details}
\label{appendsec:bugpattern}
%这里我们给出我们使用爬虫爬取的bug pattern的统计学数据信息
%我们爬虫了多少 github issue
We implemented a crawler to collect all \textbf{816} GitHub issues from the four code repositories. We then filtered the issues and identified \textbf{34} issues related to consensus protocol logic bugs, which were used to construct our bug patterns to help us integrate expert knowledge. Examples are shown below:

\PatternBox{Bug Pattern Details Show Case 1}
{
\textbf{Pattern ID:} Byzantine DoS via insufficient QC validation - single malicious node can stall consensus by forging high-view QCs \\
\textbf{Protocol Type:} BFT (HotStuff Variant) \\
\textbf{Category:} Security / Liveness / Insufficient Validation

\vspace{0.5em}
\textbf{Description:}
A critical vulnerability exists in the \texttt{AggregateQC} creation and verification process. A single Byzantine node can indefinitely stall consensus by providing a forged Quorum Certificate (QC) with a deceptively high view number.

\vspace{0.5em}
\textbf{Root Cause Analysis:}
\begin{itemize}
    \item \texttt{CreateAggregateQC} blindly trusts and collects QCs from \texttt{TimeoutMsg} payloads without individual verification.
    \item \texttt{VerifyAggregateQC} only validates the \texttt{highQC} (the QC with the highest view), creating a single point of failure.
\end{itemize}

\vspace{0.5em}
\textbf{Attack Scenario:}
\begin{enumerate}
    \item \textbf{Forgery:} A Byzantine node crafts a \texttt{TimeoutMsg} containing a forged QC with a very high view number.
    \item \textbf{Aggregation:} The Leader trustingly includes this invalid \texttt{highQC} into the \texttt{AggregateQC}.
    \item \textbf{Proposal:} The Leader proposes the new block/view containing this \texttt{AggregateQC}.
    \item \textbf{Rejection:} Honest nodes attempt to verify the \texttt{AggregateQC}. Since \texttt{VerifyQuorumCert(highQC)} fails, the entire proposal is rejected.
    \item \textbf{Stall:} As the leader cannot form a valid proposal that honest nodes accept, the consensus process halts.
\end{enumerate}

\vspace{0.5em}
\textbf{Impact:}
Complete blockchain liveness failure. A single malicious actor can perpetually halt the network's progress.

\vspace{0.5em}
\textbf{Fix \& Mitigation:}
\begin{itemize}
    \item \textbf{Individual Verification:} Modify \texttt{CreateAggregateQC} to independently validate every QC using \texttt{VerifyQuorumCert} before its inclusion.
    \item \textbf{Zero-Trust Policy:} Never trust aggregate structures from other nodes without verifying each component.
\end{itemize}

\vspace{0.5em}
\textbf{Context:}
\begin{itemize}
    \item \textbf{Reference:} \href{https://github.com/relab/hotstuff/issues/224}{relab/hotstuff Issue \#224}
    \item \textbf{Vulnerable Functions:} \texttt{CreateAggregateQC}, \texttt{VerifyAggregateQC}
    \item \textbf{Confirmed by:} meling (maintainer)
\end{itemize}
}

\PatternBox{Bug Pattern Details Show Case 2}
{
\textbf{Pattern ID:} VerifyLeader OOM - notify list grows unbounded when futures not removed from repl.notify on error \\
\textbf{Protocol Type:} CFT (Hashicorp Raft) \\
\textbf{Bug Category:} Liveness / Memory Leak / OOM

\vspace{0.5em}
\textbf{Description:}
In the \texttt{VerifyLeader} mechanism, a node failure triggers an unbounded memory growth. The \texttt{verifyLeader()} function registers a verify future \texttt{v} into the \texttt{repl.notify} list for each replica. When a heartbeat or \texttt{AppendEntries} returns an error, the cleanup process is incomplete, leading to a memory leak that eventually causes an Out-Of-Memory (OOM) crash.

\vspace{0.5em}
\textbf{Root Cause Analysis:}
\begin{itemize}
    \item \textbf{Asymmetric Cleanup:} The future \texttt{v} is tracked in two distinct structures: \texttt{leaderState.notify} and \texttt{repl.notify}.
    \item \textbf{Incomplete Removal:} On error, the future is only deleted from \texttt{leaderState.notify} but remains in \texttt{repl.notify}.
    \item \textbf{Resource Exhaustion:} The replication goroutine persists in holding references to failed or completed futures, preventing Garbage Collection (GC) and causing the list to grow infinitely.
\end{itemize}

\vspace{0.5em}
\textbf{Lessons \& Best Practices:}
\begin{itemize}
    \item \textbf{Symmetric Lifecycle:} Ensure that any future or callback registered in multiple structures is removed from \textit{all} of them upon completion or failure.
    \item \textbf{Failure Path Testing:} Monitor memory usage specifically under failure conditions, as leaks often remain hidden during healthy cluster operations.
    \item \textbf{Cleanup Patterns:} Utilize \texttt{defer} or RAII-like patterns to guarantee cleanup across all possible execution paths.
    \item \textbf{Subscription Management:} For async notifications, consider explicit lifecycle management or weak references to manage subscription longevity.
\end{itemize}

\vspace{0.5em}
\textbf{Context:}
\begin{itemize}
    \item \textbf{Reference:} \href{https://github.com/hashicorp/raft/issues/291}{Hashicorp Raft Issue \#291}
    \item \textbf{Tags:} \texttt{cft}, \texttt{memory\_leak}, \texttt{verify\_leader}, \texttt{async\_notify}
\end{itemize}
}

\section{Constraints Details}
\label{appendsec:constraints}
%这里我们给出我们constraints的例子。
%分析bft cft不同

Here, we provide examples of constraints generated for BFT and CFT systems, which help generate scenarios tailored to different types of consensus protocols.

%BFT Constraint
\PatternBox{Protocol Constraints \& BFT Attack Model}
{
\textbf{Protocol Type Constraint:} BFT (Byzantine Fault Tolerant) \\
This model assumes an environment capable of tolerating \textbf{Byzantine faults}, where nodes may behave maliciously or arbitrarily.

\vspace{0.8em}
\textbf{I. CFT-type Attack Vectors (Crash Faults):}
\begin{itemize}
    \item \textbf{Availability:} Node crash, recovery, and frequent restarts.
    \item \textbf{Network:} Partitions, significant latency, packet loss, and message reordering.
    \item \textbf{Logic:} Boundary condition exploits and duplicate message utilization.
\end{itemize}

\vspace{0.5em}
\textbf{II. BFT-specific Attack Vectors (Malicious Behavior):}
\begin{itemize}
    \item \textbf{Equivocation:} Sending conflicting votes/proposals for the same round/height to different peers.
    \item \textbf{Payload Tampering:} Injecting invalid, malformed, or tampered protocol messages.
    \item \textbf{Selective Broadcast:} Distributing messages to a subset of nodes to induce state inconsistency.
    \item \textbf{Strategic Withholding:} A malicious leader intentionally delaying or withholding proposals to stall liveness.
    \item \textbf{State Forgery:} Crafting fraudulent certificates, votes, or cryptographic signatures.
    \item \textbf{Timing Attacks:} Exploiting vulnerabilities within the protocol's timing and timeout assumptions.
\end{itemize}

\vspace{0.8em}
\textbf{III. Attack Mindset \& Objectives:}
\textit{Assuming control of up to $f$ malicious nodes ($f < n/3$), the attacker aims to:}
\begin{enumerate}
    \item \textbf{Safety Violation:} Can honest nodes be coerced into accepting an invalid state?
    \item \textbf{Liveness Violation:} Can the system be prevented from reaching consensus indefinitely?
    \item \textbf{Consistency Violation:} Can different honest nodes be led to perceive divergent ``truths''?
\end{enumerate}

\vspace{0.5em}
\textbf{Note:} Effective testing must simulate a clever attacker who strategically exploits the protocol's specific logic weaknesses rather than just random fault injection.
}

%CFT Constraint
\PatternBox{Protocol Constraints \& CFT Fault Model}
{
\textbf{Protocol Type Constraint:} CFT (Crash Fault Tolerant) \\
This model assumes an environment that can only tolerate \textbf{crash faults}. All nodes are assumed to be honest; they either follow the protocol correctly or stop functioning entirely.

\vspace{0.8em}
\textbf{I. Allowed Attack Methods (Non-Malicious Faults):}
\begin{itemize}
    \item \textbf{Node Dynamics:} Random node crashes (immediate stop) and restarts (recovery with potential loss of in-memory state).
    \item \textbf{Network Adversity:} Network partitions, extreme communication delays, message loss, and packet reordering.
    \item \textbf{Implementation Limits:} Boundary conditions (overflows, null values, extreme inputs) and duplicate resource utilization.
\end{itemize}

\vspace{0.5em}
\textbf{II. Prohibited Attack Methods (Byzantine Behavior):}
\begin{itemize}
    \item \textbf{No Malicious Messages:} Sending forged, tampered, or intentionally malformed payloads is \textit{strictly prohibited}.
    \item \textbf{No Equivocation:} Double voting or sending conflicting proposals in the same round is not allowed.
    \item \textbf{No Selective Broadcast:} Intentional inconsistency through partial message distribution is excluded.
    \item \textbf{No Identity Theft:} Actions involving signature forgery or impersonating other nodes are disallowed.
\end{itemize}

\vspace{0.8em}
\textbf{III. Core Assumptions:}
\textit{In a CFT environment, the testing focus shifts from "malice" to "reliability":}
\begin{enumerate}
    \item \textbf{Binary State:} Nodes are either "Correct" (following the protocol) or "Failed" (stopped). There is no "half-dead" or Byzantine behavior.
    \item \textbf{Honest Participation:} All participating nodes are assumed to act in good faith to reach a common consensus.
    \item \textbf{Recovery Consistency:} The system must maintain safety and eventually regain liveness after network partitions heal or nodes recover.
\end{enumerate}

\vspace{0.5em}
\textbf{Note:} Testing for CFT protocols should exhaustively simulate unstable infrastructure and edge-case logic rather than adversarial intent.
}



\section{Memory Details}
\label{appendsec:memory}

% repository knowledge and bug patterns 给个例子
% 大概几kb

The memory contains repository knowledge and bug patterns, with a small total size of about 360 KB. Examples of repository knowledge are shown below.

\PatternBox{Repository Knowledge: etcd Raft (Go)}
{
\textbf{Core Testing Helpers \& API:}
\begin{itemize}
    \item \texttt{newNetwork(peers...)}: Creates an in-memory Raft network simulation environment, supporting fault injection and message dropping.
    \item \texttt{newTestRaft(id, election, heartbeat, storage)}: Initializes a Raft instance for white-box testing with specific timing parameters.
    \item \texttt{nextEnts(r, s)}: Returns entries that are committed but not yet applied, and advances the \texttt{applied} index.
    \item \texttt{MemoryStorage}: In-memory implementation of the \texttt{Storage} interface. Key operations like \texttt{Append}, \texttt{CreateSnapshot}, and \texttt{Compact} must respect boundary conditions (e.g., \texttt{Compact} cannot exceed \texttt{lastIndex}).
\end{itemize}

\vspace{0.5em}
\textbf{Key States \& Message Types:}
\begin{itemize}
    \item \textbf{States:} \texttt{StateFollower} (0), \texttt{StateCandidate} (1), \texttt{StateLeader} (2), \texttt{StatePreCandidate} (3).
    \item \textbf{Local Triggers:} \texttt{MsgHup} (triggers election), \texttt{MsgBeat} (triggers heartbeat broadcast).
    \item \textbf{Replication:} \texttt{MsgProp} (proposal), \texttt{MsgApp} (append entries), \texttt{MsgAppResp} (append response).
    \item \textbf{Control:} \texttt{MsgHeartbeat}, \texttt{MsgSnap} (snapshot transfer), \texttt{MsgVote} (solicit vote).
\end{itemize}

\vspace{0.8em}
\textbf{Lessons Learned \& Common Pitfalls:}
\begin{itemize}
    \item \textbf{Implicit No-op Entry:} Upon becoming a Leader (\texttt{becomeLeader}), the node automatically appends an empty \textbf{no-op} entry. This must be accounted for when calculating expected \texttt{lastIndex}.
    \item \textbf{Single-Node Quorum:} In a 1-node cluster, entries are committed immediately upon append, but still require a call to \texttt{nextEnts} to drive the application logic.
    \item \textbf{Flow Control (Inflights):} When \texttt{MaxInflight=1}, the system is highly sensitive. Probe responses (empty \texttt{MsgApp}) must be handled correctly to unpause the pipeline and prevent \texttt{PipelineStall}.
    \item \textbf{Snapshot Boundaries:} \texttt{MemoryStorage} maintains a dummy entry at \texttt{FirstIndex()-1}. If \texttt{HardState.Commit} is less than \texttt{FirstIndex}, Raft will panic during initialization.
    \item \textbf{Learner Promotion Safety:} Promoting a Learner to a Voter while it has unacknowledged inflight messages can lead to state machine stalls (e.g., stuck in \texttt{StateProbe}) if the transition logic does not explicitly clear buffers.
\end{itemize}

\vspace{0.5em}
\textbf{Execution Context:} 
Always run tests using the package context (\texttt{go test -v -run TestName .}) rather than standalone files (\texttt{go test file.go}) to ensure that unexported helper functions in \texttt{raft\_test.go} are correctly linked.
}

\begin{table}[!t]
\caption{Codebase information.}
\vspace{0.5em}
\label{tab:codebase}
\centering
\resizebox{0.7\textwidth}{!}{
\begin{tabular}{l c l l r r }
\toprule
\textbf{Repo} & \textbf{Type} & \textbf{Protocol} & \textbf{Lang} & \textbf{Stars} & \textbf{Forks} \\
\midrule
etcd~\cite{etcdioraft} & CFT & Raft & Go & 982 & 222 \\
efficientEPaxos~\cite{efficientepaxos} & CFT & EPaxos & Go & 631 & 138 \\
relabHotStuff~\cite{relabhotstuff} & BFT & HotStuff & Go & 217 & 67 \\
sui~\cite{suimystenlabs} & BFT & BullShark & Rust & 7.6k & 11.7k \\
\bottomrule
\end{tabular}
}
\end{table}




\section{Codebase Details}
\label{appendsec:codebase}
%appendix: 代码仓star很多 
% 多种不同的变成语言  go, Rust %reseach 还是production
We evaluate our approach on four popular consensus protocol codebases, including both research prototypes and production-level repositories, as summarized in Table~\ref{tab:codebase}.


% \begin{table}[!t]
% % 1. 设置较小的全局列间距
% \setlength{\tabcolsep}{3pt} 
% % 2. 使用 @{\hspace{宽度}} 精确控制间距
% \begin{tabularx}{\columnwidth}{@{} l @{\hspace{4pt}} c @{\hspace{4pt}} l @{\hspace{4pt}} c @{\hspace{4pt}} r @{\hspace{4pt}} r @{} } 
% \toprule
% \textbf{Repo} & \textbf{Type} & \textbf{Protocol} & \textbf{Lang} & \textbf{Stars} & \textbf{Forks}\\
% \midrule
% % etcd 行  star 982 fork 222
% {etcd-\ref{tab:codebase}} & CFT & 
% \href{https://raft.github.io/raft.pdf}{Raft} & Go & 982 & 222 \\
% \midrule
% % efficient/epaxos 行  star 631 fork 138
% \href{https://github.com/efficient/epaxos}{efficient} & CFT & 
% \href{https://dl.acm.org/doi/pdf/10.1145/2517349.2517350}{EPaxos} & Go & 631 & 138 \\
% \midrule
% % relab/hotstuff 行  star 217 fork 67
% \href{https://github.com/relab/hotstuff}{relab} & BFT & 
% \href{https://dl.acm.org/doi/pdf/10.1145/3293611.3331591}{HotStuff} & Go & 217 & 67 \\
% \midrule
% % sui/bullshark 行 star 7.6k fork 11.7k
% \href{https://github.com/MystenLabs/sui}{sui} & BFT & 
% \href{https://dl.acm.org/doi/pdf/10.1145/3548606.3559361}{Narwhal/Bullshark} & Rust & 7.6k & 11.7k \\
% \bottomrule
% \end{tabularx}
% \end{table}



% Figure~\ref{fig:totaltokens} 显示了我们ablation studies中total tokens使用的情况,
% 虽然 去掉 bug-exploitation 和reflection-loop可以轻微减少token usage但是去很难发现有用的bugs,去掉其他的componnent都导致消耗更多的tokens,尤其是去掉scenario-generator,会回退到baseliens.
% 结果证明了我们的模块可以帮助\consensustester 消耗更少的tokens发现更多有意义的bugs.


% bug-exploitation 
%  state-analyzer
%  constraints-analyzer
%  scenario-generator
%  reflection-loop



\section{Two Types of Bug Example}
\label{appendsec:bugexample}
\subsection{Implementation Bug}
\label{appendsec:implementationbug}
% https://docs.qq.com/doc/DQU9Mc3hrRXRtZHJn?no_promotion=1 第一个

\BugAnalysisBox{Implementation Bug}
{
\textbf{Incorrect RPC Message Dispatch (Copy-Paste Error)}

\paragraph{Problem Description}
A critical logic error exists in the \texttt{bcastTryPreAccept} function. Despite correctly populating a \texttt{TryPreAccept} structure (\texttt{tpa}), the code erroneously references an unrelated \texttt{PreAccept} structure (\texttt{pa}) during dispatch.

\paragraph{Affected Components}
\begin{itemize}
    \item \textbf{File:} \texttt{src/epaxos/epaxos.go}
    \item \textbf{Function:} \texttt{bcastTryPreAccept} (Lines 580--595)
\end{itemize}

\paragraph{Root Cause Analysis}
In \texttt{epaxos.go}, the function spends lines 588--594 filling the fields of \texttt{tpa}. However, at line 595, the code executes: \texttt{args := \&pa}. Since \texttt{pa} represents a different RPC schema, the broadcast sends stale or misaligned protocol data.


\paragraph{Attack Scenario / Impact}
\begin{enumerate}
    \item \textbf{Mangled Communication:} The leader broadcasts a message intended as \texttt{TryPreAccept}, but the payload is cast from a \texttt{PreAccept} structure.
    \item \textbf{Protocol Collapse:} 
        \begin{itemize}
            \item \textbf{Recovery Failure:} Doubtful instances remain unresolved, leading to a permanent Liveness hang.
            \item \textbf{State Machine Divergence:} Peers may accept inconsistent dependency graphs based on outdated \texttt{pa} data.
        \end{itemize}
\end{enumerate}

\paragraph{Corrective Action}
Update the reference at line 595 to the correct initialized structure:
\begin{center}
    \texttt{Change \textbf{\&pa} $\rightarrow$ \textbf{\&tpa}}
\end{center}
} % <--- 确保这个右大括号存在，且前面没有未闭合的环境

\subsection{Logic Bug}
\label{appendsec:logicbug}

\subsubsection{BFT Bug}
\label{appendsec:bftbug}
% hotstuff 4
% bft one
\BugAnalysisBox{BFT Bug}
{
\textbf{Case Study: Safety Violation via Cascading Forks in HotStuff}

\paragraph{Affected Components}
\begin{itemize}
    \item \textbf{Files:} \texttt{protocol/rules/chainedhotstuff.go}, \texttt{protocol/consensus/voter.go}
    \item \textbf{Modules:} \texttt{ChainedHotStuff.VoteRule()}, \texttt{CommitRule()}
\end{itemize}

\paragraph{Root Cause Analysis}
The vulnerability stems from an interplay between weak liveness checks and the 3-chain commitment logic:
\begin{itemize}
    \item \textbf{Vulnerable Vote Rule:} The liveness condition (\texttt{QC.View > bLock.View}) is too permissive. If a node's \texttt{bLock} is lagged (e.g., after a reboot), any Quorum Certificate (QC) from a higher view can bypass the safety check, even if it extends an ancient branch.
    \item \textbf{Context-Free Commit Rule:} The \texttt{CommitRule} strictly validates the 3-chain structure (consecutive views and matching parent hashes) but fails to ensure the chain is part of the "canonical" branch.
    \item \textbf{Mutable bLock:} \texttt{bLock} updates whenever a 2-chain is formed on \textit{any} fork, lowering the barrier for Byzantine leaders to push a malicious branch.
\end{itemize}

\paragraph{Attack Scenario (n=4, f=1)}
\begin{enumerate}
    \item \textbf{State:} Honest nodes hold a high QC for $B_3$ (Chain: $G \rightarrow B_1 \rightarrow B_2 \rightarrow B_3$). A target honest node has \texttt{bLock = Genesis}.
    \item \textbf{Step 1 (View 4):} A Byzantine leader proposes $B_4'$ extending $B_1$. Since $QC_{B1}.View (1) > bLock.View (0)$, the node votes, forming $QC_4'$.
    \item \textbf{Step 2 (Views 5--6):} Consecutive Byzantine leaders extend the fork to $B_5'$ and $B_6'$. The node's \texttt{bLock} is updated to $B_4'$, facilitating the fork's progression.
    \item \textbf{Step 3 (View 7):} Upon receiving $B_7'$, the \texttt{CommitRule} detects a valid 3-chain ($B_4' \rightarrow B_5' \rightarrow B_6'$). Consequently, $B_4'$ is committed, bypassing the correct $B_2 \rightarrow B_3$ sequence.
\end{enumerate}

\paragraph{Impact \& Security Analysis}
This attack is critical as it leads to a \textbf{Safety Violation}. Unlike transient liveness issues, this results in permanent state pollution and ledger divergence across honest nodes. The attack is feasible in scenarios involving node restarts or state synchronization lags, especially when the protocol allows a Byzantine node to be elected as a leader for consecutive views.

\paragraph{Verification}
The exploit was successfully reproduced using:\\
\texttt{go test -v -run TestByzantineLeaderBuildsCascadingForkAcrossViews}
}


\subsubsection{CFT Bug}
\label{appendsec:cftbug}
%cft一个
%epaxos 2

\BugAnalysisBox{CFT Bug}
{
\textbf{Case Study: Dependency Divergence during EPaxos Recovery}

\paragraph{Affected Components}
\begin{itemize}
    \item \textbf{File:} \texttt{src/epaxos/epaxos.go}
    \item \textbf{Modules:} \texttt{handlePrepareReply()}, \texttt{startPhase1()}, and \texttt{updateAttributes()}
\end{itemize}

\paragraph{Root Cause Analysis}
The vulnerability is rooted in the improper re-computation of command attributes during the recovery phase. When a replica attempts to recover an instance whose Fast Path status is ambiguous, it incorrectly invokes \texttt{startPhase1}, which triggers a fresh call to \texttt{updateAttributes}.
\begin{itemize}
    \item \textbf{Attribute Volatility:} Instead of preserving the original dependencies recorded during the pre-accept phase, \texttt{startPhase1} re-scans the local log for conflicting commands.
    \item \textbf{Dependency Smuggling:} If a new, conflicting command has arrived during the leader's downtime, it will be "smuggled" into the recovered command's dependency set. This creates a discrepancy between the dependencies used during the original Fast Path execution and the final committed state.
\end{itemize}

\paragraph{Attack Scenario (3 Nodes: $R_0, R_1, R_2$)}
\begin{enumerate}
    \item \textbf{Fast Path Execution:} $R_0$ proposes $Cmd_A$ (Key="x"). $R_0$ and $R_1$ accept it with $\text{deps}=\{\}$. $R_0$ reaches a Fast Quorum, commits, and executes $Cmd_A$ locally ($x=1$).
    \item \textbf{Leader Crash:} $R_0$ crashes immediately before broadcasting the \texttt{Commit} message to others.
    \item \textbf{Conflict Injection:} $R_2$ proposes $Cmd_B$ (conflicting with $Cmd_A$). $R_2$ and $R_1$ accept $Cmd_B$. Crucially, $R_1$ now sees $Cmd_B$ in its log.
    \item \textbf{Vulnerability Trigger:} $R_1$ initiates recovery for $Cmd_A$. Failing to observe a Fast Quorum, it executes \texttt{startPhase1}. The \texttt{updateAttributes} function sees $Cmd_B$ and re-calculates $Cmd_A$'s dependencies as $\text{deps}=\{Cmd_B\}$.
    \item \textbf{State Divergence:} The cluster eventually commits $Cmd_A$ with $\text{deps}=\{Cmd_B\}$, implying an execution order of $B \rightarrow A$. However, the failed leader $R_0$ has already executed $A$ with $\text{deps}=\{\}$.
\end{enumerate}

\paragraph{Impact \& Security Analysis}
This bug represents a \textbf{Linearizability Violation}. Upon recovery, $R_0$ will find its local execution history contradicts the globally committed serialized order. In high-throughput systems with frequent leader transitions or network partitions, this divergence leads to permanent database corruption and "split-brain" states at the application level.

\paragraph{Verification}
The divergence can be reproduced by a controlled test case where:
\begin{itemize}
    \item A specific node is isolated after a Fast Path commit.
    \item A conflicting command is injected into the remaining quorum.
    \item Recovery is forced on the surviving nodes.
\end{itemize}
}

\subsubsection{Bug Triggers from Bug Exploitation}
\label{appendsec:deepbug}
%bug exploritation的一个
%分析
%hotstuff 5
\BugAnalysisBox{Bug from Bug Exploitation}
{
\textbf{Case Study: History Erasure via Deep Fork from Genesis}

\paragraph{Affected Components}
\begin{itemize}
    \item \textbf{Files:} \texttt{protocol/rules/chainedhotstuff.go}, \texttt{protocol/consensus/voter.go}
    \item \textbf{Modules:} \texttt{ChainedHotStuff.VoteRule()}, \texttt{NewChainedHotStuff()}
\end{itemize}

\paragraph{Root Cause Analysis}
This vulnerability represents the most extreme manifestation of "HighQC Ignorance." It exploits the discrepancy between the \texttt{HighQC} (the highest known certificate) and \texttt{bLock} (the locked block) during edge cases such as node recovery.
\begin{itemize}
    \item \textbf{Lock Amnesia:} The current implementation stores \texttt{bLock} in memory without persistence. Upon a node restart or crash-recovery, \texttt{bLock} is reset to \texttt{Genesis} (View 0).
    \item \textbf{Trivial Safety Satisfaction:} The Safety condition in \texttt{VoteRule} checks if a proposal extends the \texttt{bLock}. When \texttt{bLock = Genesis}, the \texttt{blockchain.Extends(block, Genesis)} call returns \texttt{true} for \textit{any} valid block in the system, effectively neutralizing the safety guardrail.
    \item \textbf{Liveness Bypass:} While the Liveness condition (\texttt{qcBlock.View > bLock.View}) typically prevents old forks, a leader can satisfy it using any QC with $\text{View} > 0$ (e.g., $QC_{B1}$) while simultaneously targeting nodes that have just reset their locks.
\end{itemize}

\paragraph{Attack Scenario: Full History Erasure}
Assume an established chain: $G \rightarrow B_1 \rightarrow B_2 \rightarrow B_3 \rightarrow B_4 \rightarrow B_5$ with $\text{HighQC} = QC_{B5}$.
\begin{enumerate}
    \item \textbf{Target Condition:} A subset of honest nodes restart, causing their \texttt{bLock} to revert to \texttt{Genesis}.
    \item \textbf{The Deep Fork:} A Byzantine leader ignores the entire history and proposes $B_6'$ extending directly from \texttt{Genesis} (Parent: \texttt{Genesis}, QC: $QC_{G}$).
    \item \textbf{Validation Failure:}
        \begin{itemize}
            \item \textbf{Liveness:} If using $QC_{B1}$, $QC_{B1}.View(1) > bLock.View(0)$, condition passes.
            \item \textbf{Safety:} \texttt{Extends($B_6'$, Genesis)} is \texttt{true}, condition passes.
        \end{itemize}
    \item \textbf{Outcome:} Honest nodes vote for $B_6'$, enabling the Byzantine leader to build an alternative chain ($G \rightarrow B_6' \rightarrow B_7' \dots$) that competes with the legitimate history, potentially leading to a massive state roll-back if the fork gains quorum.
\end{enumerate}

\paragraph{Impact \& Security Analysis}
This bug exposes the system to \textbf{Total History Erasure}. Unlike standard forks that deviate by a few blocks, this "Deep Fork" allows an attacker to rewrite the ledger from the Genesis block. While HotStuff's voting rules (one vote per view) prevent simultaneous double-commits within the same view, the "Lock Amnesia" creates a window where a Byzantine leader can force honest nodes to abandon the canonical chain in favor of a malicious alternative, causing severe consistency issues and service disruption.

\paragraph{Verification}
Reproduced via: \\
\texttt{go test -v -run TestByzantineLeaderBuildsDeepForkFromGenesis} \\
\textit{Result:}  \texttt{BUG CONFIRMED: Voter accepted proposal extending Genesis despite HighQC being at View 5.}
}

This bug is uncovered through further exploration of the bug described below.

\BugAnalysisBox{Bug: Design Vulnerability in Consensus Voting Logic}
{
\textbf{Case Study: Byzantine Leader Proposes Block Ignoring HighQC}

\paragraph{Affected Components}
\begin{itemize}
    \item \textbf{Files:} \texttt{protocol/rules/chainedhotstuff.go}, \texttt{protocol/consensus/voter.go}
    \item \textbf{Modules:} \texttt{ChainedHotStuff.VoteRule()}, \texttt{NewChainedHotStuff()}
\end{itemize}

\paragraph{Root Cause Analysis}
This issue stems from a design-implementation gap where \texttt{VoteRule} validates proposals against \texttt{bLock} (Locked Block) rather than the \texttt{HighQC} (Highest Quorum Certificate). Under specific edge cases—such as node recovery or state synchronization—\texttt{bLock} can become desynchronized from the actual chain height.

\begin{itemize}
    \item \textbf{Implementation Logic:} Per the HotStuff paper, a vote is cast if:
    \begin{enumerate}
        \item \textbf{Liveness:} \texttt{QC.View > bLock.View} 
        \item \textbf{Safety:} The block extends \texttt{bLock}.
    \end{enumerate}
    \item \textbf{Persistence Failure:} \texttt{bLock} is initialized to \texttt{Genesis} and kept only in memory. It is updated during the \texttt{CommitRule} but is not persisted to disk.
    \item \textbf{The Vulnerability:} Upon restart, a node's \texttt{bLock} reverts to \texttt{Genesis} (View 0). Even if the blockchain is recovered, the node will accept any proposal that satisfies the Liveness condition against a view of 0, even if the proposal ignores the actual canonical tip (\texttt{HighQC}).
\end{itemize}

\paragraph{Attack Scenarios}
\textbf{Scenario: Node Recovery / New Joiner}
\begin{enumerate}
    \item \textbf{State:} Chain exists as $G \rightarrow B_1 \rightarrow B_2 \rightarrow B_3$ (\texttt{HighQC} = $B_3$).
    \item \textbf{Trigger:} A node restarts; \texttt{bLock} resets to \texttt{Genesis}.
    \item \textbf{The Attack:} A Byzantine leader proposes $B_4'$ extending $B_1$ using $QC_{B1}$.
    \item \textbf{Validation:} \texttt{VoteRule} sees $QC_{B1}.View(1) > bLock.View(0)$. The Liveness condition passes, and the honest node mistakenly votes for a fork that ignores $B_3$.
\end{enumerate}

\paragraph{Verification}
Reproduced via: \texttt{go test -v -run TestByzantineLeaderIgnoresHighQC ./protocol/consensus/}
\begin{quote}
\texttt{ BUG DETECTED: Voter.Verify() accepted malicious proposal extending B1 instead of HighQC B3} \\
\texttt{ BUG CONFIRMED: Honest node voted for malicious proposal!} \\
\texttt{Result: This allows a Byzantine leader to fork the chain and split quorum.}
\end{quote}

\paragraph{Impact \& Security Analysis}
While this vulnerability strictly follows the HotStuff paper's safety rules (preventing double-commits within the same view via \texttt{lastVotedView} checks), it creates significant \textbf{Liveness and Stability risks}:
\begin{itemize}
    \item \textbf{Chain Forking:} Malicious leaders can force honest nodes to vote on stale branches.
    \item \textbf{Quorum Fragmentation:} Voting power is split between the canonical chain and the fork, stalling finality.
    \item \textbf{Increased Latency:} Recovery from such "HighQC ignorance" requires extra views to re-synchronize the locked state.
\end{itemize}
}


\section{Bug Example for Five Classes}
\label{appendsec:bugclassexample}
% Recovery\&Execution Divergence
% % epxaos 1
\BugAnalysisBox{Recovery\&Execution Divergence}
{
\textbf{CommittedCommandVanishing (Loss of Committed State)}

\paragraph{Affected Components}
\begin{itemize}
    \item \textbf{File:} \texttt{src/epaxos/epaxos.go}
    \item \textbf{Line Numbers:} 1542--1556
    \item \textbf{Logic:} The \texttt{else} branch in the recovery handler where the leader receives no information from the quorum.
\end{itemize}

\paragraph{Root Cause Analysis}
The vulnerability arises because the recovery leader assumes an instance is "empty" or "unused" if it receives no records from its recovery quorum. When this happens, it erroneously proposes a \texttt{NO-OP} to finalize the instance. This logic fails to account for \textbf{Amnesia} (node restarts losing volatile state) or \textbf{Network Partitions} that isolate the only nodes holding the original \texttt{Commit} record.

\paragraph{Attack Scenario (5 Nodes: $R_0, R_1, R_2, R_3, R_4$)}
\begin{enumerate}
    \item \textbf{Fast Path Execution:} $R_0$ proposes $Cmd_A$ (\texttt{Key="k", Val=999}). Nodes $R_0, R_1, R_2$ accept it (Fast Quorum).
    \item \textbf{Partial Commit:} $R_0$ sends a \texttt{Commit} message \textbf{only} to $R_2$. $R_2$ marks the command as \texttt{Committed}.
    \item \textbf{Compound Failure:} 
        \begin{itemize}
            \item $R_0$ crashes permanently.
            \item $R_1$ crashes and reboots (Amnesia), wiping its log.
            \item $R_2$ is partitioned from the rest of the cluster.
        \end{itemize}
    \item \textbf{Flawed Recovery:} $R_3$ initiates recovery for the same instance and queries $\{R_1, R_3, R_4\}$. 
    \item \textbf{Voiding the Instance:} $R_1, R_3,$ and $R_4$ all report no record. $R_3$ incorrectly concludes the instance was never used and proposes \texttt{NO-OP}.
    \item \textbf{State Divergence:} $R_1, R_3, R_4$ commit and execute \texttt{NO-OP}. When the partition heals, $R_2$ has already executed $Cmd_A$, causing a permanent state divergence.
\end{enumerate}

\paragraph{Impact \& Risk Assessment}
\begin{itemize}
    \item \textbf{Outcome:} Violation of the "Stability" property and Linearizability.
    \item \textbf{Feasibility:} High in environments without strict persistence (Sync-to-Disk) and frequent network partitions.
    \item \textbf{Risk Level:} \textbf{Medium} (Requires specific failure interleaving).
\end{itemize}

\paragraph{Corrective Action}
The protocol should either:
\begin{itemize}
    \item Strictly prohibit Amnesia by requiring persistent WAL (Write-Ahead Logs).
    \item Enhance the recovery logic to require an explicit "Proof of Non-Existence" before proposing a \texttt{NO-OP}.
\end{itemize}
}

% Persistence\&Monotonicity Violations
% % epaoxs 8
\BugAnalysisBox{Persistence\& Monotonicity Violations}
{
\textbf{Case Study: ExecutedStateAmnesia (Double Execution Vulnerability)}

\paragraph{Problem Description}
In the \texttt{efficient-epaxos} implementation, the \texttt{EXECUTED} status of an instance and the execution progress tracker (\texttt{ExecedUpTo}) are not persisted to stable storage. Upon a replica crash and recovery, the system loses its execution cursor, leading to the re-application of already committed commands to the state machine.

\paragraph{Affected Components}
\begin{itemize}
    \item \textbf{Files:} \texttt{src/epaxos/epaxos.go}, \texttt{src/epaxos/epaxos-exec.go}
    \item \textbf{Variables:} \texttt{ExecedUpTo} (Execution cursor), \texttt{inst.Status}
\end{itemize}

\paragraph{Root Cause Analysis}
The vulnerability stems from the volatility of the execution state:
\begin{enumerate}
    \item \textbf{Volatile Progress Tracking:} The \texttt{ExecedUpTo} array is stored purely in memory. On reboot, it resets to $-1$, losing the boundary between executed and pending commands.
    \item \textbf{Status Volatility:} The transition to \texttt{EXECUTED} (\texttt{epaxos-exec.go:128}) is never recorded via \texttt{recordInstanceMetadata}.
    \item \textbf{Asynchronous Gap:} Since execution happens asynchronously after the \texttt{COMMITTED} phase, the metadata update logic (which usually triggers on commit) fails to capture the final "executed" lifecycle of the instance.
\end{enumerate}



\paragraph{Attack Scenario (Double Execution)}
\begin{enumerate}
    \item \textbf{Initial State:} $R_0$ commits and executes $(0,0): \text{INC}(x)$, $(0,1): \text{INC}(x)$. State $x=2$, \texttt{ExecedUpTo[0]} = 1.
    \item \textbf{Crash \& Amnesia:} $R_0$ crashes. \texttt{ExecedUpTo} is wiped. 
    \item \textbf{Recovery:} $R_0$ reboots, restores instances $(0,0)$ and $(0,1)$ as \texttt{COMMITTED} from disk, but its memory-based cursor is $-1$.
    \item \textbf{Re-execution:} The \texttt{executeCommands} loop re-scans from index 0 ($ExecedUpTo + 1$) and applies the commands again.
    \item \textbf{Divergence:} Final state $x=4$, while non-crashing replicas remain at $x=2$.
\end{enumerate}

\paragraph{Impact \& Security Analysis}
\begin{itemize}
    \item \textbf{Outcome:} Violation of \textbf{Exactly-Once Semantics}.
    \item \textbf{Risk Level:} \textbf{Critical} for financial ledgers or non-idempotent state transitions.
    \item \textbf{Integrity:} Causes irreversible state divergence that the consensus protocol cannot automatically repair.
\end{itemize}

\paragraph{Verification}
Reproduced via: \texttt{go test -v src/epaxos/execedupto\_amnesia\_test.go} \\
\textit{Result:} \texttt{FAIL: State machine mismatch. Expected x=3, got x=6.}
}


% Dependency\&Topology Flaws
% %epxos 4


\BugAnalysisBox{Dependency\&Topology Flaws}
{
\textbf{DependencyDivergence (Execution Order Inversion)}

\paragraph{Problem Description}
The recovery mechanism may fail to preserve established dependency relationships. If the "bridge node"—the only node that witnessed the causal link between two commands—undergoes amnesia (state loss), the recovery process may treat ordered commands as concurrent. This leads to an execution order inversion through deterministic tie-breaking, causing state machine divergence.

\paragraph{Affected Components}
\begin{itemize}
    \item \textbf{File:} \texttt{src/epaxos/epaxos.go}
    \item \textbf{Line Numbers:} 1471--1488
    \item \textbf{Logic:} The \texttt{handlePrepareReply} logic where dependencies are merged or updated during the recovery phase based on quorum responses.
\end{itemize}

\paragraph{Root Cause Analysis}
EPaxos relies on nodes recording dependencies ($deps$) to serialize conflicting commands. The bug occurs when the only evidence of a dependency (e.g., $Cmd_B$ depends on $Cmd_A$) exists on a node that subsequently fails and loses its memory-based log. During recovery, the remaining nodes—each only aware of one of the commands—reconstruct the dependency graph without the causal link, mistakenly concluding the commands are concurrent.

\paragraph{Attack Scenario (5 Nodes: $R_0, R_1, R_2, R_3, R_4$)}
\begin{enumerate}
    \item \textbf{Split Proposal:} $R_1$ proposes $Cmd_A$; $R_0$ proposes $Cmd_B$ (both target Key $x$).
    \item \textbf{Bridge Node Setup:} 
        \begin{itemize}
            \item $R_2$ processes $Cmd_A$ first, then $Cmd_B$. It records $Cmd_B.deps = \{Cmd_A\}$.
            \item $R_3$ only sees $PreAccept(Cmd_A)$; $R_4$ only sees $PreAccept(Cmd_B)$. Both record empty dependencies.
        \end{itemize}
    \item \textbf{Initial Execution:} $R_2$ receives \texttt{Commit} for both and executes $A \rightarrow B$. Final state at $R_2$: $x=2$.
    \item \textbf{The Wipe (Amnesia):} $R_0, R_1$ crash permanently. $R_2$ crashes and reboots with an empty log (Amnesia).
    \item \textbf{Flawed Recovery:} 
        \begin{itemize}
            \item $R_3$ recovers $Cmd_A$. Quorum $\{R_3, R_4, R_2\}$ reports no dependency for $A$ on $B$.
            \item $R_4$ recovers $Cmd_B$. Quorum $\{R_4, R_3, R_2\}$ reports no dependency for $B$ on $A$.
            \item The link $B \rightarrow A$ is lost. Both are committed with $deps = \{\}$.
        \end{itemize}
    \item \textbf{Execution Inversion:} The surviving nodes apply tie-breaking ($R_0.id < R_1.id$). $Cmd_B$ is executed before $Cmd_A$. Final state: $x=1$.
\end{enumerate}

\paragraph{Impact \& Risk Assessment}
\begin{itemize}
    \item \textbf{Outcome:} \textbf{Linearizability Violation}. The system history is rewritten, reversing the order of operations already visible to the user.
    \item \textbf{Risk Level:} \textbf{High} in high-concurrency environments with non-persistent deployments.
    \item \textbf{Feasibility:} Highly likely if network latency is asymmetric and nodes do not use synchronous disk writes for every \texttt{PreAccept}.
\end{itemize}

\paragraph{Corrective Action}
\begin{itemize}
    \item \textbf{Persistence:} Enforce stable storage for all metadata changes to prevent state loss upon reboot.
    \item \textbf{Recovery Logic:} Modify the recovery leader's criteria to require a "Strong Quorum" or "Epoch-based" validation to ensure that any previously committed dependency link is not ignored simply because a node is empty.
\end{itemize}
}


% % Message Binding\&Signature Violationsa
% % %hotstuff 8
\BugAnalysisBox{Message Binding \& Signature Violations}
{
\textbf{Case Study: HighQC View Mismatch (View Tampering via UpdateHighQC)}

\paragraph{Problem Description}
The \texttt{UpdateHighQC} function checks if a block exists and if its view is higher than the current \texttt{HighQC}. However, it \textbf{fails to verify that the View stored within the QC matches the View of the block it certifies}. This allows an attacker to inject a QC with a tampered (artificially inflated) View field into the node's state.

\paragraph{Affected Components}
\begin{itemize}
    \item \textbf{File:} \texttt{protocol/viewstates.go}
    \item \textbf{Function:} \texttt{UpdateHighQC}
    \item \textbf{Dependencies:} \texttt{hotstuff.QuorumCert}, \texttt{blockchain} module
\end{itemize}

\paragraph{Root Cause Analysis}
The implementation logic compares the block's inherent view (\texttt{newBlock.View()}) for the threshold check but then blindly persists the entire \texttt{qc} object without consistency validation. 
\begin{itemize}
    \item \textbf{Consistency Gap:} There is no assertion that \texttt{qc.View() == newBlock.View()}.
    \item \textbf{Implicit Trust:} The function assumes the \texttt{QuorumCert} provided is internally consistent, even though the \texttt{View} field in the QC is not cryptographically bound to the signature in this specific implementation.
\end{itemize}



\paragraph{Attack Scenario (View Inflation)}
\begin{enumerate}
    \item \textbf{Setup:} The network is currently at \texttt{View v}.
    \item \textbf{Acquisition:} A Byzantine node obtains a valid \texttt{QC} certifying \texttt{Block\_A} from \texttt{View v-k}.
    \item \textbf{Tampering:} The attacker constructs a \texttt{fakeQC} with the original valid signature for \texttt{Block\_A} but an inflated \texttt{View: v + 1000}.
    \item \textbf{Replay:} The attacker sends this \texttt{fakeQC} to an honest node.
    \item \textbf{State Corruption:} Because \texttt{Block\_A.View()} exists and is valid, the node updates its \texttt{highQC} to \texttt{fakeQC}. The node now erroneously believes the highest certified view is \texttt{v + 1000}, causing it to reject legitimate messages from the actual current view.
\end{enumerate}

\paragraph{Impact \& Security Analysis}
\begin{itemize}
    \item \textbf{Liveness Violation:} Disrupts View Synchronization, locking honest nodes out of consensus.
    \item \textbf{Category:} \textbf{CWE-20: Improper Input Validation} (Cross-field consistency check missing).
    \item \textbf{Severity:} \textbf{High}. State-pollution that can be triggered remotely.
\end{itemize}

\paragraph{Corrective Action}
Add an explicit consistency check between the QC and the block metadata:
\begin{quote}
\texttt{if qc.View() != newBlock.View() \{ \\
\quad return false, fmt.Errorf("QC view mismatch") \\
\}}
\end{quote}
}



% Resource\&Operational Visibility Violations
% % Sui 1

\BugAnalysisBox{Resource \& Operational Visibility Violations}
{
\textbf{Case Study: BlockManager Memory Exhaustion (OOM via Unbounded MissingBlocks)}

\paragraph{Problem Description}
The \texttt{BlockManager} maintains a \texttt{missing\_blocks} set to track parent blocks that have not yet been received. This data structure lacks an upper bound on its capacity, allowing a Byzantine validator to flood the system with blocks referencing non-existent ancestors. This leads to an Out-of-Memory (OOM) condition and potential Denial of Service (DoS).

\paragraph{Affected Components}
\begin{itemize}
    \item \textbf{Module:} \texttt{BlockManager}
    \item \textbf{Data Structure:} \texttt{missing\_blocks} (Set/Map)
    \item \textbf{Source Note:} Explicitly acknowledged by a \texttt{TODO} comment regarding Byzantine causal history and OOM risks.
\end{itemize}



\paragraph{Vulnerability Details}
\begin{enumerate}
    \item \textbf{Unbounded Growth:} There is no hard limit on the number of entries in the \texttt{missing\_blocks} collection. Attackers can continuously produce blocks with fake causal histories.
    \item \textbf{Future-Flooding Attack:} A naive "lowest round first" eviction strategy can be bypassed. Attackers can use high-round garbage blocks to force the eviction of legitimate, low-round missing blocks that the node actually needs.
\end{enumerate}

\paragraph{Attack Scenario (Future-Flooding)}
\begin{enumerate}
    \item \textbf{Normal State:} An honest node is at \texttt{Round 10} and is missing a legitimate parent block $P$ (from \texttt{Round 10}).
    \item \textbf{Flooding:} A Byzantine attacker floods the network with blocks referencing thousands of fake parent blocks at \texttt{Round 1000}.
    \item \textbf{Malicious Eviction:} If a capacity limit exists but uses a simple round-based eviction, the node will discard the legitimate \texttt{Round 10} entry to make room for the "newer" \texttt{Round 1000} entries.
    \item \textbf{Liveness Failure:} The honest node can never complete the causal history for \texttt{Round 10}, stalling its consensus progress.
\end{enumerate}

\paragraph{Impact \& Security Analysis}
\begin{itemize}
    \item \textbf{Vulnerability Type:} \textbf{CWE-770: Allocation of Resources Without Limits}.
    \item \textbf{Attack Cost:} Extremely low. A single committee member can trigger this with low-bandwidth flooding over time.
    \item \textbf{Consequence:} System-wide crash (OOM) or permanent stalling of honest nodes.
\end{itemize}

\paragraph{Corrective Action}
\begin{itemize}
    \item \textbf{Bound the Collection:} Implement a strict maximum size for \texttt{missing\_blocks}.
    \item \textbf{Smart Eviction:} Use a \textbf{Round-Proximity} strategy. Prioritize missing blocks closer to the current round and implement per-validator quotas.
\end{itemize}
}

\subsection{Final Summary}

\section{Prompt Example}
\label{appendsec:prompt}
%是否进一步拆成5个componnet的

\subsection{\agenta Prompt}
%分析是两个subtasks


\begin{agentbox}[Orchestrator Agent]{
    % --- System ---
    \textbf{System:}
    \begin{Verbatim}[fontsize=\scriptsize]
You are the Orchestrator Agent, the central coordinator of a multi-agent system for discovering

vulnerabilities in distributed consensus protocols.

Your responsibilities:
1. Analyze confirmed bugs discovered by the Strategy Agent and TestGen Agent
2. Generate comprehensive bug reports with root cause analysis
3. Extract reusable vulnerability patterns for testing other protocol implementations
4. Track resource usage (tokens, cost) for each bug discovery

Be technical, precise, and actionable. Focus on providing insights that help developers understand

and fix the vulnerabilities.
    \end{Verbatim}
    
    % --- User ---
    \textbf{User:}
    
    \begin{Verbatim}[fontsize=\scriptsize]
As the Orchestrator Agent, analyze the following confirmed vulnerability and generate a comprehensive report.

## Confirmed Bug Details

The Strategy Agent generated an attack scenario, and the TestGen Agent confirmed it as a real vulnerability.

**Protocol**: {protocol_name}
**Protocol Type**: {protocol_type.upper()}
**Bug Name**: {scenario_name}
**Target Component**: {target_component or 'N/A'}
**Attack Category**: {attack_category or 'N/A'}

**Vulnerability Hypothesis**:
{vulnerability_hypothesis or 'N/A'}

**Preconditions**:
{preconditions_text}

**Attack Steps**:
{steps_text}

**Expected Bug Behavior**: {expected_bug_behavior or 'N/A'}
**Correct Behavior**: {correct_behavior or 'N/A'}

**Test File**: {test_file}
**Test Output** (excerpt):
```
{test_output[:2000] if test_output else 'N/A'}
```

## Resource Usage for This Bug Discovery

- **Prompt Tokens**: {prompt_tokens:,}
- **Completion Tokens**: {completion_tokens:,}
- **Total Tokens**: {total_tokens:,}
- **API Cost**: ${api_cost:.4f}

---

## Your Task

As the Orchestrator Agent, generate:

1. **Bug Report**: A comprehensive analysis including:
   - Root cause of the vulnerability
   - Potential security impact
   - Recommended fixes

2. **Bug Pattern**: A reusable pattern for testing similar vulnerabilities in other consensus protocol implementations.

## Output Format

<bug_report>
[Your detailed bug report with root cause analysis, impact assessment, and \\
fix recommendations]
</bug_report>

<bug_pattern>
Pattern Name: [A concise, descriptive name for this vulnerability pattern]
Bug Category: [safety/liveness/agreement]
Fault Type: [e.g., byzantine_leader, network_partition, crash_recovery]
Trigger Condition: [Specific conditions that trigger this vulnerability]
Description: [Detailed description of the vulnerability pattern and why it  occurs]
Test Template: [How to test for this pattern in other protocol implementations]
</bug_pattern>
    \end{Verbatim}
    
    % --- Assistant ---
    \bigskip
    }
    \end{agentbox}

\subsection{\agentb Prompt}
% 2个componnet
\begin{agentbox}[Strategy Agent]{
    % --- System ---
    \textbf{System:}
    \begin{Verbatim}[fontsize=\scriptsize]
You are a distributed systems security researcher, specializing in discovering vulnerabilities in consensus

protocol implementations.

## Your Goal

For the target protocol, devise **cunning, creative, ORIGINAL** test scenarios to discover:
- **Safety Violations**: Different nodes commit different values for the same slot
- **Liveness Violations**: System cannot make progress
- **Agreement Violations**: Nodes disagree on the committed state

## Available Tool

You have access to ONE tool:

### get_pattern_details(pattern_id: str)
Query detailed information about a known bug pattern. Use this when you see a pattern_id 
in the bug pattern list that seems relevant and you want to learn more before deciding.

**Usage Guidelines**:
- Only call this for patterns you're genuinely considering using
- The pattern_id is a summary - if it already seems irrelevant, don't query it
- If the tool warns that a pattern was already used, you MUST create a substantially different attack vector

## Knowledge Base

Use your knowledge of:
- Distributed systems theory
- Consensus protocol design (Paxos, Raft, PBFT, HotStuff, Tendermint, etc.)
- Common implementation pitfalls
- Byzantine fault tolerance
- Network partitions and timing attacks
- State machine replication
- Crash recovery scenarios

## Repository Context (If Provided)

- Use it to ground your scenario in realistic components and constraints.
- Prefer components or behaviors that are likely to exist in the target repo.

## Bug Pattern Context (If Provided)

- You may receive a list of pattern IDs (which summarize the bug).
- Use `get_pattern_details(pattern_id)` to learn more about any interesting pattern.
- Select **at most one** relevant pattern and adapt it to the target protocol.
- **DO NOT reuse a pattern that has already been used** - the tool will warn you if it was.
- If no pattern fits or all have been used, create an **original** scenario.

## Thinking Approach

1. **Understand the Protocol Type**
   - What is the fault model? (CFT vs BFT)
   - What are the core safety/liveness assumptions?
   - What are the key invariants that must hold?

2. **Identify Attack Vectors**
   Think about common vulnerability categories:
   - Race conditions in concurrent message handling
   - State persistence and recovery bugs
   - View change / leader election edge cases
   - Quorum intersection violations
   - Timeout and timing-related issues
   - Message ordering and delivery assumptions
   - Equivocation and double-voting scenarios
   - Lock state management errors

3. **Creative Scenario Design**
   - Consider edge cases that developers might miss
   - Think about interactions between multiple components
   - Imagine adversarial network conditions
   - Consider crash-restart scenarios at critical moments

4. **Design Verification**
   - How do we determine if the attack succeeded?
   - What specific assertions can capture this issue?

## IMPORTANT: Avoid Unrealistic Scenarios

Before proposing any attack scenario, verify it is REALISTIC:

1. Can the trigger condition actually happen in a real deployment?
2. Are there upstream checks that prevent the scenario?
3. Is the fault model (CFT vs BFT) correctly applied?

Examples of UNREALISTIC scenarios to avoid:
- Testing CFT protocol with Byzantine behavior (CFT assumes honest nodes)
- Assuming attacker can bypass all input validation
- Assuming network delivers messages in impossible orders
- Creating test states that cannot exist in normal operation

**A scenario is only valid if it can genuinely occur in production!**

## Output Format

<thinking>
[Understanding the Protocol]
...I observe that the key mechanism of this protocol is...
...It assumes...
...The fault model is...

[Identifying Attack Vectors]
...Common vulnerabilities in this type of protocol include...
...I should consider...

[Creative Reasoning]
...If I construct a scenario where...
...This might lead to...
...The key insight is...

[Designing Verification]
...If the attack succeeds, I should observe...
...The assertion should check...
</thinking>

<attack_scenario>
Name: [Descriptive name, e.g., "CrashInducedLockAmnesiaLeadingToConflictingCommits"]

Target Component: [Specific component/function being attacked]

Vulnerability Hypothesis: [What problem you believe exists]

Attack Category: [safety / liveness / agreement]

Inspiration Source: [pattern_memory::<pattern_id> | original]

Preconditions:
1. [What initial state is needed]
2. [...]

Attack Steps:
1. [Specific action]
2. [...]
3. [...]

Expected Bug Behavior: [What will happen if there's a bug]

Correct Behavior: [How the protocol should respond]

Assertions:
1. [How to verify the attack succeeded]
2. [...]
</attack_scenario>
    \end{Verbatim}
    
    % --- User ---
    \textbf{User:}
    
    \begin{Verbatim}[fontsize=\scriptsize]
## Task

Devise a **cunning, creative, ORIGINAL** attack scenario for {protocol_name}.

Use your knowledge of distributed systems and consensus protocols to identify potential vulnerabilities.
Think about:
- What could go wrong in the implementation?
- What edge cases might developers miss?
- How could an adversary (or faulty node) exploit timing or state management?

**Remember**: 
- Ensure your scenario is REALISTIC and can genuinely occur in production!
- Be CREATIVE and think of scenarios that might not be obvious
- Focus on implementation-level bugs, not protocol design flaws

Now think deeply about potential vulnerabilities and design an attack scenario...
    \end{Verbatim}
    \textbf{Bug Exploitation}
    
    \begin{Verbatim}[fontsize=\scriptsize]
You are a distributed systems security researcher in BUG EXPLOITATION MODE.

A bug has been CONFIRMED in the target protocol. Your mission is to leverage this 
discovery to uncover MORE vulnerabilities by exploring related attack vectors.

## Exploitation Strategies

When exploiting a confirmed bug, consider these approaches:

### 1. Root Cause Analysis
- What is the fundamental assumption being violated?
- Where else in the codebase might this assumption be violated?
- Are there other code paths that share the same flawed logic?

### 2. Variant Generation
- Can the same bug be triggered through different entry points?
- What happens with different timing? Different message orders?
- Are there edge cases the original test didn't cover?

### 3. Pattern Propagation
- If this component has this bug, do similar components have it too?
- Does the bug exist in related protocols (view change, leader election, etc.)?
- Can the bug be combined with other faults to cause more severe issues?

### 4. Deeper Exploration
- What happens if we push the bug further?
- Can we escalate from liveness to safety violations?
- Are there cascading effects in dependent components?

## Key Principles

1. **Build on the Discovery**: Use the confirmed bug as a foundation
2. **Think Creatively**: The bug reveals a blind spot - where else might it exist?
3. **Avoid Repetition**: Generate genuinely NEW scenarios, not minor variations
4. **Stay Realistic**: All scenarios must be achievable in real deployments

## Output Format

<thinking>
[Analyzing the Confirmed Bug]
...The root cause of this bug is...
...This reveals an assumption that...

[Identifying Related Attack Vectors]
...Based on this, I should explore...
...A similar vulnerability might exist in...

[Designing New Scenario]
...If I target X component with Y approach...
...This is different from the original because...
</thinking>

<attack_scenario>
Name: [New descriptive name - must be different from original]
Target Component: [May be same or different component]
Vulnerability Hypothesis: [Related but distinct vulnerability]
Attack Category: [safety / liveness / agreement]
Relation to Original Bug: [How this relates to the confirmed bug]

Preconditions:
1. [Required conditions]
2. [...]

Attack Steps:
1. [Specific action]
2. [...]

Expected Bug Behavior: [What happens if this variant exists]
Correct Behavior: [Expected correct protocol response]

Assertions:
1. [Verification method]
2. [...]
</attack_scenario>
    \end{Verbatim}
    
    % --- Assistant ---
    \bigskip
    }
\end{agentbox}
\bigskip
\subsection{\agentc Prompt}
%分析是三个subtasks
\begin{agentbox}[TestGen Agent]{
    % --- System ---
    \textbf{System:}
    \begin{Verbatim}[fontsize=\scriptsize]
You are a test code generation expert, skilled at transforming attack scenarios into executable test code.

## Your Goal

Transform the attack scenario generated by Strategy Agent into:
1. **Repository-style compliant** test code
2. **Executable, verifiable** test cases
3. **Clear assertions** to determine if the attack succeeded

## Available Tools for Repository Knowledge

You have these powerful tools to leverage accumulated knowledge:

### `repo_knowledge(compact=True)`
Get cached knowledge about the repository. Use `compact=True` (default) for LLM-optimized summary.
- Contains: test structure, coding style, helper functions with signatures and usage examples
- Contains: lessons learned from previous errors

### `repo_knowledge_find_lessons(error_message)`
**CALL THIS FIRST** when encountering compilation or test errors.
- Searches previous lessons that match the error
- Returns solutions that worked before

### `repo_knowledge_add_lesson(issue, solution, error_pattern="")`
**ALWAYS CALL THIS** after successfully fixing an error.
- Records the issue and solution for future reference
- `error_pattern`: Optional regex to match similar errors

### `repo_knowledge_add_helper(name, file, purpose, signature="", usage_example="", returns="")`
Call this when you discover a useful helper function.
- Records detailed information including signature and usage example
- Helps future test generation use the helper correctly

## Workflow

1. **Analyze Repository Structure**
   - Call `repo_knowledge()` to get test structure, coding style, and helper functions
   - Look for helpers with `signature` and `usage_example` - use them directly
   - Check `lessons_learned` for common pitfalls to avoid
   - Find similar test files as reference if needed

2. **Generate Test Code**
   - Follow the repository's coding style exactly
   - Use helper functions with their documented signatures
   - Use correct package/module declarations
   - Import necessary dependencies

3. **Add Assertions**
   - Detect expected vulnerability behavior
   - If protocol is correct, test should PASS
   - If there's a bug, test should FAIL

4. **Execute Test**
   - Write the test file
   - Run the test command
   - Collect output results

## Code Quality Requirements

- **Correct package name**: Read from existing files, don't guess
- **Correct imports**: Use modules that actually exist in the repository
- **Follow style**: Mimic how other tests in the repository are written
- **Clear comments**: Explain what each step is doing

## Error Handling Workflow (CRITICAL)

When test compilation/execution fails, follow this **exact** sequence:

### Step 1: Search Previous Lessons FIRST
```
error_message = "..." # the error you got
-> Call repo_knowledge_find_lessons(error_message)
```
If a relevant lesson is found, apply the solution directly.

### Step 2: Investigate If No Lesson Found
- Use Bash to run commands: `ls`, `grep`, etc.
- Read relevant files to understand the issue
- Check if imports, package names, paths are correct

### Step 3: Fix and Verify
- Make corrections based on your investigation
- Run the test again

### Step 4: RECORD THE LESSON (Don't forget!)
After successfully fixing:
```
-> Call repo_knowledge_add_lesson(
    issue="What went wrong (be specific)",
    solution="How you fixed it (be detailed)",
    error_pattern="optional regex to match similar errors"
)
```

### Step 5: Update Helper Details If Discovered
If you found a helper function's signature or usage during investigation:
{helper_example}

**Key Principle**: Learn from errors and RECORD what you learn. 
Every error you fix should benefit future test generation.

## Output Format

<analysis>
[Analyzing repository test structure]
...repo_knowledge shows test files are in...
...Naming convention is...
{analysis_example}
...Previous lessons: avoid X, do Y instead...
</analysis>

<test_code language="{lang_order}">
// Complete test code (see language-specific templates below)
</test_code>

<execution_plan>
File Path: [Path where test file will be written]
Test Command: [Command to run the test]
</execution_plan>
    \end{Verbatim}
    
    % --- User ---
    \textbf{User:}
    
    \begin{Verbatim}[fontsize=\scriptsize]
## Language-Specific Template

{language_template}
"""
    )

    prompt_parts.append(
        f"""
## Target Protocol

- **Name**: {protocol_name}
- **Language**: {language}

## Attack Scenario

{attack_scenario}
"""
    )

    if repo_knowledge:
        prompt_parts.append(
            f"""
## Repository Knowledge (Pre-loaded)

{repo_knowledge}
"""
        )

    # Add test commands
    test_commands = get_test_commands(language)
    prompt_parts.append(test_commands)

    prompt_parts.append(
        """
## Task

1. Review the repo_knowledge above (especially helper signatures and lessons learned)
2. Use the language-specific template as a guide for code structure
3. Transform the attack scenario into test code using available helpers
4. Ensure code can compile and run
5. Execute test and report results
6. If errors occur: search lessons -> fix -> record lesson

Now begin...
    \end{Verbatim}
    
    % --- Assistant ---
    \bigskip
    \textbf{Go Test Template}
    
    \begin{Verbatim}[fontsize=\scriptsize]
## Go Test Template

```go
package <package_name>

import (
    "testing"
    "time"
    // Add other imports from repo_knowledge
)

// Test<ScenarioName> tests <brief description>
//
// Attack scenario: <scenario description>
// Expected behavior: <what the test verifies>
func Test<ScenarioName>(t *testing.T) {
    // Step 1: Setup cluster/nodes
    // Use helpers from repo_knowledge (e.g., NewNetwork, MakeCluster)
    
    // Step 2: Establish initial state
    // - Elect leader if needed
    // - Submit baseline transactions
    
    // Step 3: Execute attack scenario
    // - Create partition / inject faults
    // - Trigger specific protocol conditions
    
    // Step 4: Verify safety/liveness properties
    // - Check for safety violations
    // - Verify expected behavior
    
    // Step 5: Cleanup and final assertions
    if safetyViolation {
        t.Fatalf("SAFETY VIOLATION: %s", details)
    }
}
```

### Go Test Commands
- Single test: `go test -v -run TestName ./path/...`
- With timeout: `go test -v -timeout 60s -run TestName ./...`
- With race detector: `go test -v -race -run TestName ./...`
    \end{Verbatim}
    \textbf{Rust Test Template}
    
    \begin{Verbatim}[fontsize=\scriptsize]
## Rust Test Template

```rust
#[cfg(test)]
mod tests {
    use super::*;
    // Add imports from repo_knowledge
    
    /// Creates a logger for test output
    fn default_logger() -> Logger {
        let decorator = slog_term::TermDecorator::new().build();
        let drain = slog_term::FullFormat::new(decorator).build().fuse();
        let drain = slog_async::Async::new(drain).build().fuse();
        Logger::root(drain, o!())
    }
    
    /// Test<scenario_name> tests <brief description>
    ///
    /// Attack scenario: <scenario description>
    /// Expected behavior: <what the test verifies>
    #[test]
    fn test_<scenario_name>() {
        let logger = default_logger();
        
        // Step 1: Setup cluster/nodes
        // Use helpers from repo_knowledge (e.g., Network::new, RawNode::new)
        
        // Step 2: Establish initial state
        // - Trigger elections
        // - Propose baseline entries
        
        // Step 3: Execute attack scenario
        // - Simulate partitions
        // - Inject faults
        
        // Step 4: Verify safety/liveness properties
        assert!(condition, "SAFETY VIOLATION: {}", details);
        
        // Step 5: Cleanup and final assertions
        println!("Test passed - no safety violations detected");
    }
}
```

### Rust Test Commands
- Single test: `cargo test test_name -- --nocapture`
- In harness: `cargo test -p harness test_name -- --nocapture`
- With logging: `RUST_LOG=debug cargo test test_name -- --nocapture`
    \end{Verbatim}
    
    \textbf{Java Test Template}
    
    \begin{Verbatim}[fontsize=\scriptsize]
## Java Test Template (JUnit 5)

```java
package <package.name>;

import org.junit.jupiter.api.Test;
import org.junit.jupiter.api.BeforeEach;
import org.junit.jupiter.api.AfterEach;
import static org.junit.jupiter.api.Assertions.*;
// Add other imports from repo_knowledge

/**
 * Test<ScenarioName> - <brief description>
 *
 * Attack scenario: <scenario description>
 * Expected behavior: <what the test verifies>
 */
public class <ScenarioName>Test {

    // Test fixtures
    private TestCluster cluster;
    
    @BeforeEach
    public void setUp() throws Exception {
        // Initialize test cluster
        // Use helpers from repo_knowledge
    }
    
    @AfterEach
    public void tearDown() throws Exception {
        // Cleanup resources
        if (cluster != null) {
            cluster.shutdown();
        }
    }
    
    @Test
    public void test<ScenarioName>() throws Exception {
        // Step 1: Setup
        // Create cluster with required configuration
        
        // Step 2: Establish initial state
        // - Wait for leader election
        // - Submit baseline transactions
        
        // Step 3: Execute attack scenario
        // - Create partition / inject faults
        // - Submit conflicting transactions
        
        // Step 4: Verify safety properties
        assertFalse(safetyViolation, "SAFETY VIOLATION: " + details);
        
        // Step 5: Heal partition and verify convergence
        assertEquals(expectedState, actualState, "States should converge after partition heals");
    }
}
```

### Java Test Commands
- Maven: `mvn test -Dtest=TestClassName`
- Maven single method: `mvn test -Dtest=TestClassName#testMethodName`
- Gradle: `gradle test --tests TestClassName`
\end{Verbatim}
}
\end{agentbox}



\end{document}


% This document was modified from the file originally made available by
% Pat Langley and Andrea Danyluk for ICML-2K. This version was created
% by Iain Murray in 2018, and modified by Alexandre Bouchard in
% 2019 and 2021 and by Csaba Szepesvari, Gang Niu and Sivan Sabato in 2022.
% Modified again in 2023 and 2024 by Sivan Sabato and Jonathan Scarlett.
% Previous contributors include Dan Roy, Lise Getoor and Tobias
% Scheffer, which was slightly modified from the 2010 version by
% Thorsten Joachims & Johannes Fuernkranz, slightly modified from the
% 2009 version by Kiri Wagstaff and Sam Roweis's 2008 version, which is
% slightly modified from Prasad Tadepalli's 2007 version which is a
% lightly changed version of the previous year's version by Andrew
% Moore, which was in turn edited from those of Kristian Kersting and
% Codrina Lauth. Alex Smola contributed to the algorithmic style files.
